\documentclass{article}
\usepackage[margin=0.8in]{geometry}
\usepackage{amsmath,amssymb,amsthm,mathtools}
\usepackage{mathrsfs,bm,bbm}
\usepackage{graphicx,booktabs,array,multirow,tabularx,longtable}
\usepackage{enumitem}
\usepackage{xcolor}
\usepackage{algorithm,algpseudocode}
\usepackage{subcaption,tikz}
\usepackage{siunitx,cancel,accents}
\usepackage{microtype}
\usepackage[numbers,sort&compress]{natbib}
\usepackage[colorlinks=true,linkcolor=blue,citecolor=blue,urlcolor=blue]{hyperref}
\usepackage{cleveref}
\setlength{\emergencystretch}{2em}
\newtheorem{assumption}{Assumption}
\newtheorem{definition}{Definition}
\newtheorem{theorem}{Theorem}
\newtheorem{lemma}{Lemma}
\newtheorem{proposition}{Proposition}
\newtheorem{openendedquestion}{Open-ended Question}
\newtheorem{conjecture}{Conjecture}
\newcommand{\coreref}[1]{\ref{#1}}

\begin{document}

\sloppy

\section*{panel\_\allowbreak{}multispell\_\allowbreak{}hankel\_\allowbreak{}irregularity}

\noindent\textit{Specialization:} v1\par

\paragraph{TL;DR.} Continuous reversible treatment histories identify each heterogeneous lag mean while making its Moore--Penrose score bounded, regular, critical, or irregular. The proposal fixes an original-coordinate geometry signature \(\Sigma\), including a concrete semialgebraic support, the maximal-Hankel-minor ideal, coordinate alignments, active support faces, and a checkable conic normal-link certificate. Real local-zeta integrals define intrinsic target-specific tail and distance-weighted bias indices, including exact logarithmic pole orders, and any normalized subordinate resolution yields the same values. Smoothness is a primitive original-coordinate Whitney condition; a frozen atlas supplies only the partition and implementation. A total cross-fitted estimator uses measurable Euclidean anchor selectors, displayed ridge normal equations, a thresholded projector estimate of the common trend, and a finite whole-class oracle grid. It attains the augmented class-level risk upper bound \(R_{G,\Sigma,k}^{\dagger}(C_{G,\Sigma,k}^{\dagger})\); no honest-length upper bound over \(\mathcal P_\Sigma\) is proved, and comparison with the deterministic \(R_{G,\Sigma,k}^{\star}\) phase remains open. The observed-law bridge proves that the mover target equals the anchor's eligible mean on \(\mathcal P_\Sigma\), without claiming that this class is a subclass of the anchor's published potential-outcome model. The framework preserves target-inactive cancellation, passes the \(K=1,T=5\) exponent-two early-kill test, recovers scalar slow-mover and critical Hankel regimes, and separately recovers finite-support root-\(G\) behavior for the \(\texttt{dist\_lag\_het}\) consumer.

\section{Project justification}

\paragraph{Gap.} Under its full potential-outcome, sampling, moment, and design hypotheses, the anchor's Theorem 4 identifies a component-specific eligible mean through a conditional-cutoff Moore--Penrose limit, but it leaves continuous-history score tails and cutoff choice unresolved. This note first proves an observed-law and target-functional bridge on the stricter fixed-signature reduced-form class \(\mathcal P_\Sigma\); it does not claim inclusion in the anchor's potential-outcome class. Scalar slow-mover theory cannot see Hankel codimension or coordinate cancellation, while existing algebraic tube results do not provide a fixed-class, boundary-compatible, two-sided probability law with distance-weighted correction.

\paragraph{Niche.} Within the anchor's induced reduced-form experiment, a strict treatment-design specialization can fix one compact full-dimensional semialgebraic mover support, positive stayer mass, density and moment envelopes, and quantitative normal-link data. Its rank ideal, target alignments, and support faces then form a decidable original-coordinate signature before outcomes are analyzed.

\paragraph{Fill.} The project defines the indices by local-zeta poles, proves their resolution-invariant tail and selector-distance implications, and specifies a total minimax-linear procedure attaining an augmented class-level risk upper bound. Comparison with the deterministic pole-sensitive oracle, same-signature minimax lower bounds, honest-length rates, and geometry adaptation are exposed as separate open questions with concrete perturbation and selection handles.

\section{Related work}

deChaisemartinDHa2025 is the immediate model and consumer. Its Theorem 4 is stated under no anticipation, independent and identically distributed full-schedule sampling, never-treated parallel trends, a separate second-moment condition, the random-coefficient distributed-lag model, \(K<T-1\), and invertibility of \(E[\Pi(M)]\); it identifies each component-specific eligible mean by a conditional-cutoff Moore--Penrose limit. Proposition \(\text{prop:published-source-bridge}\) independently verifies the induced reduced-form requirements and target-functional equality on \(\mathcal P_\Sigma\), rather than asserting that \(\mathcal P_\Sigma\) is a subclass of the published potential-outcome model. GrahamPowell2012 informs the scalar benchmark through Proposition 1.1, Theorem 2.1, and equation (29); the uncorrected \(G^{-2/3}\) balance is recovered literally, while the distance-weighted first-order correction has \(\zeta_{\Sigma,0,1}=2\) and therefore a different \(G^{-4/5}\) balance. ConcaEtAl2018 informs the nested Hankel rank ideals and singular locus. PhongSturm2000 informs the principalization and local-zeta calculation, and Wongkew1993 and DinhPham2018 provide upper geometric envelopes rather than the proposed two-sided target-specific law. KhanTamer2010 motivates the inverse-weight rate frontier. SasakiUra2022 informs scalar derivative correction, while SasakiUra2026 informs the stayer-plus-slow-mover architecture; neither treats shared Hankel strata. Donoho1994 and CaiLow2004 inform minimax and honest-length moduli, and ArmstrongKolesar2018 informs bias-aware intervals. ArellanoBonhomme2012 anchors fixed-\(T\) random-coefficient identification. GentzkowShapiroSinkinson2011 supplies the reversible newspaper-count application used by \(\texttt{dist\_lag\_het}\).

\section{Setup}

Target (estimand / structural parameter): \(\tau_k=E_P[B_k\mid X\ne0]\); for \(P\in\mathcal P_\Sigma\), Proposition \(\text{prop:published-source-bridge}\) proves that this causal mover mean equals the published induced reduced-form eligible mean \(\overline\beta_k=E_P[B_k\mid e_k\in\operatorname{Im}(M(X)')]\).
Primary functional / estimator / diagnostic: For \(P\in\mathcal P_\Sigma\), \(\tau_k=\lim_{C\to\infty}E_P[e_k'M(X)^+(\Delta Y-\Gamma)\mathbf 1\{H_k(X)\le C\}\mid X\ne0]=\overline\beta_k\). At finite \(C\), this note's indicator/mover-denominator functional equals \(\rho_{k,C}\) times the anchor's conditional-cutoff functional, where \(\rho_{k,C}=P\{H_k(X)\le C\mid X\ne0\}\); only their limits are asserted equal. The deterministic fixed-signature oracle has the proved pole-sensitive upper scale \(R_{G,\Sigma,k}^{\star}\), while the minimax-linear estimator attains \(R_{G,\Sigma,k}^{\dagger}(C_{G,\Sigma,k}^{\dagger})\), and proving \(\Omega_{G,\Sigma,k,s}(C_{G,\Sigma,k}^{\star})\lesssim R_{G,\Sigma,k}^{\star}\) remains open.
\begin{itemize}
  \item \texttt{G} --- \texttt{Positive integer.}; \texttt{Number of independently sampled groups.}
  \item \texttt{T} --- Integer with \(T\ge 2\).; \texttt{Fixed panel length.}
  \item \texttt{K} --- \texttt{Nonnegative integer.}; \texttt{Largest carryover lag.}
  \item \texttt{p} --- \texttt{Positive integer.}
  \par\smallskip\noindent\emph{Definition.} \(p=K+1\).
  \item \texttt{r} --- \texttt{Positive integer.}
  \par\smallskip\noindent\emph{Definition.} \(r=T-K-1\).
  \item \texttt{d} --- \texttt{Positive integer.}
  \par\smallskip\noindent\emph{Definition.} \(d=T-1\).
  \item \texttt{k} --- Index in \(\{0,\ldots,K\}\).
  \item \texttt{s} --- \texttt{Positive integer.}; \texttt{Hölder order and boundary-\allowbreak{}extrapolation order.}
  \item \texttt{q} --- Real number with \(q>4\).; \texttt{Moment exponent.}
  \item \texttt{\textbackslash{}\allowbreak{}eta} --- Real number in \(0,1/2)\).; \texttt{Noncoverage probability.}
  \item \texttt{\textbackslash{}\allowbreak{}mathbf c} --- \texttt{Vector of fixed positive constants.}
  \par\smallskip\noindent\emph{Definition.} \(\mathbf c=(\underline\pi,\overline\pi,\underline f,\overline f,C_q,\underline\sigma,\overline\sigma,L_s)\).
  \item \texttt{\textbackslash{}\allowbreak{}delta\_\allowbreak{}0} --- \texttt{Positive real number.}; \texttt{Certified common radius for original-\allowbreak{}coordinate conic neighborhoods and Whitney extension.}
  \item \texttt{P} --- \texttt{Probability law.}; Laws of fixed-\(T\) panel observations.
  \item \texttt{P\textasciicircum{}\{\textbackslash{}\allowbreak{}mathrm\{fs\}\allowbreak{}\}\allowbreak{}} --- \texttt{Probability law.}; Laws of fixed-\(T\) panel observations with a finite mover support.; Separate finite-support comparison law, not a member of the absolutely continuous class \(\mathcal P_\Sigma\).
  \item \texttt{D} --- \texttt{Random vector.}; \(\mathbb R^T\).; \texttt{Treatment path.}
  \item \texttt{X} --- \texttt{Random vector.}; \(\mathbb R^d\).; \texttt{Treatment-\allowbreak{}increment history.}
  \par\smallskip\noindent\emph{Definition.} \(X_a=D_{a+1}-D_a\) for \(1\le a\le d\).
  \item \texttt{Y} --- \texttt{Random vector.}; \(\mathbb R^T\).; \texttt{Outcome path.}
  \item \texttt{\textbackslash{}\allowbreak{}Delta Y} --- \texttt{Random vector.}; \(\mathbb R^r\).
  \par\smallskip\noindent\emph{Definition.} \(\Delta Y_i=Y_{K+1+i}-Y_{K+i}\) for \(1\le i\le r\).
  \item \texttt{W\_\allowbreak{}g} --- \texttt{Observed random vector.}; \texttt{Group-\allowbreak{}level observation.}
  \par\smallskip\noindent\emph{Definition.} \(W_g=(\Delta Y_g,X_g)\), \(1\le g\le G\).
  \item \texttt{B} --- \texttt{Random vector.}; \(\mathbb R^p\).; \texttt{Time-\allowbreak{}invariant heterogeneous distributed-\allowbreak{}lag coefficients.}
  \item \texttt{\textbackslash{}\allowbreak{}varepsilon} --- \texttt{Random vector.}; \(\mathbb R^r\).; \texttt{Differenced disturbance.}
  \item \texttt{\textbackslash{}\allowbreak{}Gamma} --- \texttt{Deterministic vector.}; \(\mathbb R^r\).; \texttt{Common untreated outcome-\allowbreak{}trend vector.}
  \item \texttt{S} --- \texttt{Compact full-\allowbreak{}dimensional subset.}; \(\mathbb R^d\).; \texttt{Fixed mover support recorded in the geometry signature.}
  \item \texttt{\textbackslash{}\allowbreak{}mathbf g} --- \texttt{Finite tuple of real polynomials.}
  \par\smallskip\noindent\emph{Definition.} \(S=\{x\in\mathbb R^d:g_b(x)\ge0,\ 1\le b\le J_0\}\), with irredundant sign conditions recording every active support face.
  \item \texttt{f\_\allowbreak{}X} --- \texttt{Density.}; \(f_X:S\to[0,\infty)\).; \texttt{Mover density.}
  \item \texttt{\textbackslash{}\allowbreak{}pi\_\allowbreak{}0} --- \texttt{Probability.}; \texttt{Stayer mass.}
  \par\smallskip\noindent\emph{Definition.} \(\pi_0=P(X=0)\).
  \item \texttt{M} --- \texttt{Polynomial matrix-\allowbreak{}valued map.}; \(M:\mathbb R^d\to\mathbb R^{r\times p}\).; \texttt{Hankel history design.}
  \par\smallskip\noindent\emph{Definition.} \(M(x)_{i,k+1}=x_{K+i-k}\) for \(1\le i\le r\) and \(0\le k\le K\).
  \item \texttt{e\_\allowbreak{}k} --- \texttt{Canonical vector.}; \(\mathbb R^p\).
  \par\smallskip\noindent\emph{Definition.} \((e_k)_a=\mathbf 1\{a=k+1\}\).
  \item \texttt{\textbackslash{}\allowbreak{}Pi} --- \texttt{Projector-\allowbreak{}valued map.}; \texttt{Orthogonal projector onto the complement of the column space.}
  \par\smallskip\noindent\emph{Definition.} \(\Pi(M)=I_r-MM^+\).
  \item \texttt{\textbackslash{}\allowbreak{}mathcal I\_\allowbreak{}\{\textbackslash{}\allowbreak{}rm rk\}\allowbreak{}} --- \texttt{Polynomial ideal.}; \texttt{Maximal-\allowbreak{}Hankel-\allowbreak{}minor rank ideal.}
  \par\smallskip\noindent\emph{Definition.} \(\mathcal I_{\rm rk}=\langle\det M_I: I\subseteq\{1,\ldots,r\},\ |I|=p\rangle\), where \(M_I\) is the \(p\times p\) row submatrix of \(M\).
  \item \texttt{Q} --- \texttt{Nonnegative polynomial.}; \texttt{Squared maximal-\allowbreak{}minor aggregate.}
  \par\smallskip\noindent\emph{Definition.} \(Q(x)=\det(M(x)'M(x))=\sum_{|I|=p}\det(M_I(x))^2\).
  \item \texttt{A\_\allowbreak{}k} --- \texttt{Nonnegative polynomial.}; \texttt{Coordinate-\allowbreak{}alignment numerator.}
  \par\smallskip\noindent\emph{Definition.} \(A_k(x)=e_k'\operatorname{adj}(M(x)'M(x))e_k\).
  \item \texttt{\textbackslash{}\allowbreak{}mathcal Z} --- \texttt{Algebraic subset.}; \texttt{Hankel rank-\allowbreak{}loss locus.}
  \par\smallskip\noindent\emph{Definition.} \(\mathcal Z=V(\mathcal I_{\rm rk})=\{x:Q(x)=0\}\).
  \item \texttt{E\_\allowbreak{}k} --- Semialgebraic subset of \(S\).; Lag-\(k\) eligibility set.
  \par\smallskip\noindent\emph{Definition.} \(E_k=\{x\in S:e_k\in\operatorname{Im}(M(x)')\}\).
  \item \texttt{H\_\allowbreak{}k} --- \texttt{Globally defined nonnegative measurable function.}; \(H_k:\mathbb R^d\to[0,\infty)\).; \texttt{Target-\allowbreak{}active pseudo-\allowbreak{}inverse leverage.}
  \par\smallskip\noindent\emph{Definition.} \(H_k(x)=\|(M(x)')^+e_k\|_2\) for every \(x\), using the Moore--Penrose inverse; in particular \(H_k(0)=0\).
  \item \texttt{\textbackslash{}\allowbreak{}overline H\_\allowbreak{}\{\textbackslash{}\allowbreak{}Sigma,\allowbreak{}k\}\allowbreak{}} --- \texttt{Extended nonnegative real number.}; \texttt{Signature-\allowbreak{}specific global leverage ceiling;\allowbreak{} it is finite exactly in the bounded-\allowbreak{}score branch derived by the tail-\allowbreak{}atlas theorem.}
  \par\smallskip\noindent\emph{Definition.} \(\overline H_{\Sigma,k}=\sup_{x\in S}H_k(x)\).
  \item \texttt{\textbackslash{}\allowbreak{}mu\_\allowbreak{}k} --- \texttt{Conditional-\allowbreak{}mean function.}; \(\mu_k(x)=E_P[B_k\mid X=x]\).; \texttt{History-\allowbreak{}specific mean lag effect.}
  \item \texttt{\textbackslash{}\allowbreak{}widetilde\textbackslash{}\allowbreak{}mu\_\allowbreak{}k} --- \texttt{Original-\allowbreak{}coordinate Whitney extension.}; \(\widetilde\mu_k:N_{\delta_0}(S)\to\mathbb R\).; A representative extending \(\mu_k\) with \(C^{s-1,1}\) regularity in the original treatment-history coordinates.
  \item \texttt{\textbackslash{}\allowbreak{}tau\_\allowbreak{}k} --- \texttt{Real-\allowbreak{}valued causal estimand.}; Mean lag-\(k\) effect among movers.
  \par\smallskip\noindent\emph{Definition.} \(\tau_k=E_P[B_k\mid X\ne0]\).
  \item \texttt{\textbackslash{}\allowbreak{}Sigma} --- \texttt{Fixed original-\allowbreak{}coordinate geometry signature.}; \texttt{Support,\allowbreak{} rank-\allowbreak{}ideal,\allowbreak{} alignment,\allowbreak{} face,\allowbreak{} and normal-\allowbreak{}link data held fixed across a minimax class.}
  \par\smallskip\noindent\emph{Definition.} \(\Sigma=(T,K,S,\mathbf g,\mathcal I_{\rm rk},(A_0,\ldots,A_K),\mathfrak S,\text{normal-link certificate record})\), where \(\mathfrak S\) is the canonical original-coordinate stratification by rank depth, alignment, and active support faces, and all entries have rational or real-algebraic codes.
  \item \texttt{\textbackslash{}\allowbreak{}mathfrak C\_\allowbreak{}\textbackslash{}\allowbreak{}Sigma} --- \texttt{Finite base-\allowbreak{}space normal-\allowbreak{}link certificate.}; Original-coordinate quantitative certificate attached to \(\Sigma\).
  \item \texttt{\textbackslash{}\allowbreak{}mathscr A\_\allowbreak{}\textbackslash{}\allowbreak{}Sigma} --- \texttt{Finite certified resolution atlas with partition of unity.}; \texttt{One frozen atlas used only for Hölder norms and implementation.}
  \item \texttt{j} --- Index of a chart in \(\mathscr A_\Sigma\).
  \item \texttt{\textbackslash{}\allowbreak{}phi\_\allowbreak{}j} --- \texttt{Proper analytic chart map.}; \(\phi_j:U_j\to S\).
  \item \texttt{\textbackslash{}\allowbreak{}omega\_\allowbreak{}j} --- \texttt{Nonnegative smooth function.}
  \par\smallskip\noindent\emph{Definition.} \((\omega_j)_j\) is the deterministic partition of unity subordinate to \(\mathscr A_\Sigma\), ordered by the certificate's canonical semialgebraic code.
  \item \texttt{\textbackslash{}\allowbreak{}mathbf q\_\allowbreak{}j} --- \texttt{Vector of nonnegative integers.}; Chart valuation of \(Q\).
  \item \texttt{\textbackslash{}\allowbreak{}mathbf a\_\allowbreak{}\{kj\}\allowbreak{}} --- \texttt{Vector of nonnegative integers.}; Chart valuation of \(A_k\).
  \item \texttt{\textbackslash{}\allowbreak{}mathbf d\_\allowbreak{}j} --- \texttt{Vector of nonnegative rational numbers.}; Chart valuation of \(\operatorname{dist}(\cdot,\mathcal Z)\).
  \item \texttt{\textbackslash{}\allowbreak{}mathbf b\_\allowbreak{}j} --- \texttt{Vector of nonnegative integers.}; \texttt{Chart valuation of the Jacobian and active boundary orthant.}
  \item \texttt{Z\_\allowbreak{}k(z)} --- Extended-real local-zeta integral indexed by real \(z\).; \texttt{Mellin integral whose first positive singularity determines the leverage-\allowbreak{}tail index.}
  \item \texttt{Z\_\allowbreak{}\{k,\allowbreak{}s\}\allowbreak{}(z)} --- Extended-real distance-weighted local-zeta integral indexed by real \(z\).; \texttt{Mellin integral whose first positive singularity determines the frontier-\allowbreak{}remainder index.}
  \item \texttt{\textbackslash{}\allowbreak{}alpha\_\allowbreak{}\{\textbackslash{}\allowbreak{}Sigma,\allowbreak{}k\}\allowbreak{}} --- Element of \( (0,\infty]\).; \texttt{Resolution-\allowbreak{}invariant lag-\allowbreak{}specific leverage-\allowbreak{}tail power.}
  \item \texttt{m\_\allowbreak{}\{\textbackslash{}\allowbreak{}Sigma,\allowbreak{}k\}\allowbreak{}} --- \texttt{Nonnegative integer.}; \texttt{Resolution-\allowbreak{}invariant tail logarithmic multiplicity.}
  \item \texttt{\textbackslash{}\allowbreak{}mathcal J\_\allowbreak{}\{\textbackslash{}\allowbreak{}Sigma,\allowbreak{}k,\allowbreak{}s\}\allowbreak{}} --- Nonnegative function on \([1,\infty]\).; \texttt{Distance-\allowbreak{}weighted frontier remainder integral.}
  \item \texttt{\textbackslash{}\allowbreak{}zeta\_\allowbreak{}\{\textbackslash{}\allowbreak{}Sigma,\allowbreak{}k,\allowbreak{}s\}\allowbreak{}} --- Element of \( (0,\infty]\).; \texttt{Resolution-\allowbreak{}invariant power of the frontier remainder.}
  \item \texttt{\textbackslash{}\allowbreak{}ell\_\allowbreak{}\{\textbackslash{}\allowbreak{}Sigma,\allowbreak{}k,\allowbreak{}s\}\allowbreak{}} --- \texttt{Nonnegative integer.}; \texttt{Resolution-\allowbreak{}invariant logarithmic multiplicity of the frontier remainder.}
  \item \texttt{\textbackslash{}\allowbreak{}mathcal E\_\allowbreak{}\{\textbackslash{}\allowbreak{}Sigma,\allowbreak{}k,\allowbreak{}s,\allowbreak{}C\}\allowbreak{}} --- \texttt{Linear boundary-\allowbreak{}extrapolation operator.}; Order-\(s-1\) Taylor extension defined by the frozen atlas.
  \item \texttt{A\_\allowbreak{}C} --- Compact subset of \(S\).; Closed finite-leverage anchor set for cutoff \(C\).
  \item \texttt{A\_\allowbreak{}\{jC\}\allowbreak{}} --- Compact chart-anchor subset of \(S\).; Part of \(A_C\) meeting the support of partition weight \(\omega_j\).
  \item \texttt{p\_\allowbreak{}C} --- Borel map from \(S\) to \(S\).; Canonical Euclidean projection onto \(A_C\) when that set is nonempty.
  \item \texttt{p\_\allowbreak{}\{jC\}\allowbreak{}} --- Borel map from \(S\) to \(S\).; Canonical chart-anchor selector with fallback to \(p_C\).
  \item \texttt{x\_\allowbreak{}\textbackslash{}\allowbreak{}Sigma} --- Point of \(S\).; \texttt{Fixed signature-\allowbreak{}coded value used only to totalize empty-\allowbreak{}anchor cases.}
  \item \texttt{\textbackslash{}\allowbreak{}mathcal D\_\allowbreak{}\{\textbackslash{}\allowbreak{}Sigma,\allowbreak{}k,\allowbreak{}s\}\allowbreak{}} --- Nonnegative function on \([1,\infty]\).; \texttt{Selector-\allowbreak{}distance integral controlling Taylor extrapolation error.}
  \item \texttt{V\_\allowbreak{}\{P,\allowbreak{}k\}\allowbreak{}} --- Nonnegative function on \([1,\infty]\).; \texttt{Law-\allowbreak{}specific truncated score-\allowbreak{}variance scale.}
  \item \texttt{\textbackslash{}\allowbreak{}overline V\_\allowbreak{}\{\textbackslash{}\allowbreak{}Sigma,\allowbreak{}k\}\allowbreak{}} --- Nonnegative function on \([1,\infty]\).; \texttt{Worst-\allowbreak{}case truncated variance over the fixed-\allowbreak{}signature class.}
  \item \texttt{\textbackslash{}\allowbreak{}overline\{\textbackslash{}\allowbreak{}mathfrak B\}\allowbreak{}\_\allowbreak{}\{\textbackslash{}\allowbreak{}Sigma,\allowbreak{}k,\allowbreak{}s\}\allowbreak{}} --- Nonnegative function on \([1,\infty]\).; \texttt{Uniform post-\allowbreak{}correction bias envelope.}
  \item \texttt{\textbackslash{}\allowbreak{}mathcal G\_\allowbreak{}\{G,\allowbreak{}\textbackslash{}\allowbreak{}Sigma\}\allowbreak{}} --- Finite subset of \([1,\infty]\).; \texttt{Tie-\allowbreak{}broken geometric cutoff grid with an infinity option.}
  \item \texttt{C\_\allowbreak{}\{G,\allowbreak{}\textbackslash{}\allowbreak{}Sigma,\allowbreak{}k\}\allowbreak{}\textasciicircum{}\{\textbackslash{}\allowbreak{}star\}\allowbreak{}} --- Element of \(\mathcal G_{G,\Sigma}\).; \texttt{Deterministic whole-\allowbreak{}class oracle cutoff.}
  \item \texttt{R\_\allowbreak{}\{G,\allowbreak{}\textbackslash{}\allowbreak{}Sigma,\allowbreak{}k\}\allowbreak{}\textasciicircum{}\{\textbackslash{}\allowbreak{}star\}\allowbreak{}} --- \texttt{Positive real number.}; \texttt{Deterministic whole-\allowbreak{}class oracle squared-\allowbreak{}risk scale.}
  \item \texttt{\textbackslash{}\allowbreak{}operatorname\{Risk\}\allowbreak{}\_\allowbreak{}\{G,\allowbreak{}k\}\allowbreak{}(\textbackslash{}\allowbreak{}Sigma)} --- \texttt{Nonnegative real number.}; \texttt{Exact minimax squared risk over the fixed-\allowbreak{}signature class.}
  \item \texttt{\textbackslash{}\allowbreak{}operatorname\{Len\}\allowbreak{}\_\allowbreak{}\{G,\allowbreak{}k\}\allowbreak{}(\textbackslash{}\allowbreak{}Sigma,\allowbreak{}\textbackslash{}\allowbreak{}eta)} --- \texttt{Nonnegative real number.}; \texttt{Exact shortest uniformly honest worst-\allowbreak{}case expected length.}
  \item \texttt{\textbackslash{}\allowbreak{}widehat\textbackslash{}\allowbreak{}Gamma} --- Estimator in \(\mathbb R^r\).
  \item \texttt{\textbackslash{}\allowbreak{}widehat\textbackslash{}\allowbreak{}tau\_\allowbreak{}k(C)} --- Real-valued total estimator indexed by \(C\in[1,\infty]\).
  \item \texttt{\textbackslash{}\allowbreak{}widehat C\_\allowbreak{}k} --- Data-measurable element of \(\mathcal G_{G,\Sigma}\).; \texttt{Feasible geometry-\allowbreak{}selected cutoff.}
  \item \texttt{\textbackslash{}\allowbreak{}mathcal C\_\allowbreak{}\{G,\allowbreak{}k\}\allowbreak{}} --- Random interval in \(\mathbb R\).; \texttt{Uniformly honest lag interval.}
  \item \texttt{\textbackslash{}\allowbreak{}mathcal P\_\allowbreak{}\textbackslash{}\allowbreak{}Sigma} --- \texttt{Law class.}; \texttt{Probability laws.}; The fixed-signature class defined by \(\text{def:fixed-signature-class}\) with displayed parameters \((T,K,s,q,\mathbf c)\) suppressed.
  \item \texttt{\textbackslash{}\allowbreak{}Lambda\_\allowbreak{}\textbackslash{}\allowbreak{}Sigma} --- \texttt{Positive real number.}; \texttt{Fixed signature-\allowbreak{}coded coordinatewise cap for minimax-\allowbreak{}linear donor weights.}
  \item \texttt{c\_\allowbreak{}\textbackslash{}\allowbreak{}Sigma} --- \texttt{Positive real number.}; \texttt{Fixed signature-\allowbreak{}and-\allowbreak{}envelope variance-\allowbreak{}penalty constant in the empirical minimax-\allowbreak{}linear objective.}
  \item \texttt{\textbackslash{}\allowbreak{}Omega\_\allowbreak{}\{G,\allowbreak{}\textbackslash{}\allowbreak{}Sigma,\allowbreak{}k,\allowbreak{}s\}\allowbreak{}} --- Nonnegative function on \( [1,\infty]\).; \texttt{Whole-\allowbreak{}class expected empirical extrapolation modulus for the minimax-\allowbreak{}linear correction.}
  \item \texttt{R\_\allowbreak{}\{G,\allowbreak{}\textbackslash{}\allowbreak{}Sigma,\allowbreak{}k\}\allowbreak{}\textasciicircum{}\{\textbackslash{}\allowbreak{}dagger\}\allowbreak{}} --- Extended nonnegative function on \(\mathcal G_{G,\Sigma}\).; \texttt{Augmented feasible squared-\allowbreak{}risk criterion containing variance,\allowbreak{} deterministic bias,\allowbreak{} empirical extrapolation,\allowbreak{} and regular-\allowbreak{}nuisance terms.}
  \item \texttt{C\_\allowbreak{}\{G,\allowbreak{}\textbackslash{}\allowbreak{}Sigma,\allowbreak{}k\}\allowbreak{}\textasciicircum{}\{\textbackslash{}\allowbreak{}dagger\}\allowbreak{}} --- Element of \(\mathcal G_{G,\Sigma}\).; \texttt{Tie-\allowbreak{}broken minimizer of the augmented feasible class criterion.}
\end{itemize}

\section{Assumptions}

\begin{assumption}[ass:at-least-square]\label{ass:at-least-square}
\(T\ge 2K+2\).
\par\smallskip\noindent\emph{Novel.} This is this paper's sufficient at-least-square specialization for the full-column-rank Hankel argument. It is not attributed to de Chaisemartin and D'Haultf\oe{}uille, who use a weaker lag restriction together with a separate projector-invertibility design condition.
\end{assumption}

\begin{assumption}[ass:iid-groups]\label{ass:iid-groups}
\((W_g)_{g=1}^G\) are independent and identically distributed with law \(P\).
\par\smallskip\noindent\emph{Standard} (Independent group sampling, \cite{GrahamPowell2012}).
\end{assumption}

\begin{assumption}[ass:distributed-lag-model]\label{ass:distributed-lag-model}
\(\Delta Y=\Gamma+M(X)B+\varepsilon\).
\par\smallskip\noindent\emph{Standard} (Time-invariant random-coefficient distributed-lag model, \cite{deChaisemartinDHa2025}).
\end{assumption}

\begin{assumption}[ass:conditional-mean-zero]\label{ass:conditional-mean-zero}
\(E_P[\varepsilon\mid X]=0\) almost surely.
\par\smallskip\noindent\emph{Standard} (Conditional mean independence of differenced disturbances, \cite{deChaisemartinDHa2025}).
\end{assumption}

\begin{assumption}[ass:stayer-mass]\label{ass:stayer-mass}
\(\underline\pi\le\pi_0\le\overline\pi<1\).
\par\smallskip\noindent\emph{Standard} (Positive stayer mass bounded away from one, \cite{SasakiUra2026}).
\end{assumption}

\begin{assumption}[ass:mover-density]\label{ass:mover-density}
\(P(X\in dx\mid X\ne0)=f_X(x)\mathbf 1\{x\in S\}\,dx\).
\par\smallskip\noindent\emph{Standard} (Absolutely continuous mover design, \cite{GrahamPowell2012}).
\end{assumption}

\begin{assumption}[ass:density-envelope]\label{ass:density-envelope}
\(\underline f\le f_X(x)\le\overline f\) for Lebesgue-almost every \(x\in S\).
\par\smallskip\noindent\emph{Standard} (Two-sided bounded mover density, \cite{GrahamPowell2012}).
\end{assumption}

\begin{assumption}[ass:fixed-signature-support]\label{ass:fixed-signature-support}
The mover support, its irredundant defining inequalities, and its active face incidences are exactly the \(S\), \(\mathbf g\), and face data recorded by the fixed signature \(\Sigma\).
\par\smallskip\noindent\emph{Novel.} A treatment design fixes dose bounds and admissible reversal histories before outcomes are sampled. Boxes and polytopes used for bounded continuous doses therefore define natural fixed-signature subclasses instead of a resolution-dependent union of laws.
\end{assumption}

\begin{assumption}[ass:normal-link-certificate]\label{ass:normal-link-certificate}
The original-coordinate certificate \(\mathfrak C_\Sigma\) supplies semialgebraic conic bi-Lipschitz parametrizations with strict algebraically checkable angle and transverse-Jacobian bounds, together with certified lower and upper link-mass intervals, for every active rank-depth, target-alignment, and support-face stratum of \(\Sigma\).
\par\smallskip\noindent\emph{Novel.} The certificate records quantitative local conic regularity of a fixed treatment-history region, not a desired probability exponent. Dose boxes and polytopes used in reversible-treatment panels have explicit face cones, while an effective semialgebraic volume-bounding algorithm verifies the strict link-mass intervals without treating numerical volume inequalities as first-order polynomial formulas.
\end{assumption}

\begin{assumption}[ass:q-moments]\label{ass:q-moments}
\(E_P[\|B\|_2^q+\|\varepsilon\|_2^q]\le C_q\).
\par\smallskip\noindent\emph{Standard} (Uniform higher-moment envelope for random coefficients and disturbances, \cite{GrahamPowell2012}).
\end{assumption}

\begin{assumption}[ass:noise-variance]\label{ass:noise-variance}
\(\underline\sigma^2 I_r\preceq E_P[\varepsilon\varepsilon'\mid X=x]\preceq\overline\sigma^2 I_r\) for Lebesgue-almost every \(x\in S\).
\par\smallskip\noindent\emph{Standard} (Uniform conditional variance nondegeneracy and boundedness, \cite{GrahamPowell2012}).
\end{assumption}

\begin{assumption}[ass:frontier-holder]\label{ass:frontier-holder}
The mover-a.e. function \(\mu_k\) has a representative admitting an original-coordinate Whitney extension \(\widetilde\mu_k\in C^{s-1,1}(N_{\delta_0}(S))\) with \(\|\widetilde\mu_k\|_{C^{s-1,1}}\le L_s\).
\par\smallskip\noindent\emph{Novel.} This is a primitive smoothness condition on the conditional mean in the observed treatment-history coordinates. It is satisfied by bounded linear, polynomial, and smooth saturation response surfaces on a neighborhood of a dose box or polytope; Sasaki and Ura's scalar derivative correction motivates the order of correction but does not supply this multivariate condition.
\end{assumption}

\section{Definitions}

\begin{definition}[def:normal-link-certificate]\label{def:normal-link-certificate}
\par\smallskip\noindent\emph{Defined object.} \texttt{\textbackslash{}\allowbreak{}mathfrak C\_\allowbreak{}\textbackslash{}\allowbreak{}Sigma}
\par\smallskip\noindent\emph{Construction.} For each original-coordinate stratum \(R\in\mathfrak S\), form its base-space normal link \(L_R(x,t)=\{v\in\mathbb S^{d-1}:x+tv\in S,\ t/2\le\operatorname{dist}(x+tv,R)\le t\}\), refined by rank depth, target activity, and support-face signs. The finite certificate gives a semialgebraic conic bi-Lipschitz parametrization of every nonempty sector and explicit \(\delta_0,c_L,C_L,c_\angle,C_\angle,c_J,C_J>0\). Polynomial identities and interval bounds certify the angle, Lipschitz, and transverse normal-Jacobian inequalities on \(0<t\le\delta_0\); an effective cylindrical-decomposition integration algorithm returns rational intervals contained in \([c_L,C_L]\) for the required Hausdorff link masses. The record contains no chart label and no asserted tail power.
\par\smallskip\noindent\emph{Inputs.} \texttt{\textbackslash{}\allowbreak{}Sigma}.
\end{definition}

\begin{definition}[def:certified-atlas]\label{def:certified-atlas}
\par\smallskip\noindent\emph{Defined object.} \texttt{\textbackslash{}\allowbreak{}mathscr A\_\allowbreak{}\textbackslash{}\allowbreak{}Sigma}
\par\smallskip\noindent\emph{Construction.} Choose the lexicographically first normalized boundary-compatible principalization certified from \(\Sigma\). It principalizes the maximal-minor ideal \(\mathcal I_{\rm rk}\), every alignment polynomial \(A_k\), every active support-face ideal and inequality, and the semialgebraic graph of \(x\mapsto\operatorname{dist}(x,\mathcal Z)\). On chart \(j\), \(Q\circ\phi_j=u_{Qj}\prod_a|z_a|^{q_{ja}}\), \(A_k\circ\phi_j=u_{Akj}\prod_a|z_a|^{a_{kja}}\), \(\operatorname{dist}(\phi_j(z),\mathcal Z)=u_{Dj}\prod_a|z_a|^{d_{ja}}\), and \(|\det D\phi_j|=u_{Jj}\prod_a|z_a|^{b_{ja}}\), with certified unit bounds on the assigned orthant. Fix a smooth subordinate partition of unity \((\omega_j)_j\), refine it into disjoint canonical semialgebraic chart cells by assigning a point to the least coded chart whose cutoff is positive, and record cutoff vanishing orders on every incident face. The atlas is frozen only for evaluating the indices and implementing extrapolation; smoothness is imposed in the original coordinates.
\par\smallskip\noindent\emph{Inputs.} \texttt{\textbackslash{}\allowbreak{}Sigma}; \texttt{def:\allowbreak{}normal-\allowbreak{}link-\allowbreak{}certificate}.
\end{definition}

\begin{definition}[def:fixed-signature-class]\label{def:fixed-signature-class}
\par\smallskip\noindent\emph{Defined object.} \texttt{\textbackslash{}\allowbreak{}mathcal P\_\allowbreak{}\textbackslash{}\allowbreak{}Sigma(T,\allowbreak{}K,\allowbreak{}s,\allowbreak{}q,\allowbreak{}\textbackslash{}\allowbreak{}mathbf c)}
\par\smallskip\noindent\emph{Construction.} \(\mathcal P_\Sigma(T,K,s,q,\mathbf c)=\{P:\text{ass:at-least-square, ass:iid-groups, ass:distributed-lag-model, ass:conditional-mean-zero, ass:stayer-mass, ass:mover-density, ass:density-envelope, ass:fixed-signature-support, ass:normal-link-certificate, ass:q-moments, ass:noise-variance, and ass:frontier-holder hold for the fixed }\Sigma\}\).
\par\smallskip\noindent\emph{Member properties.} \texttt{ass:\allowbreak{}at-\allowbreak{}least-\allowbreak{}square}; \texttt{ass:\allowbreak{}iid-\allowbreak{}groups}; \texttt{ass:\allowbreak{}distributed-\allowbreak{}lag-\allowbreak{}model}; \texttt{ass:\allowbreak{}conditional-\allowbreak{}mean-\allowbreak{}zero}; \texttt{ass:\allowbreak{}stayer-\allowbreak{}mass}; \texttt{ass:\allowbreak{}mover-\allowbreak{}density}; \texttt{ass:\allowbreak{}density-\allowbreak{}envelope}; \texttt{ass:\allowbreak{}fixed-\allowbreak{}signature-\allowbreak{}support}; \texttt{ass:\allowbreak{}normal-\allowbreak{}link-\allowbreak{}certificate}; \texttt{ass:\allowbreak{}q-\allowbreak{}moments}; \texttt{ass:\allowbreak{}noise-\allowbreak{}variance}; \texttt{ass:\allowbreak{}frontier-\allowbreak{}holder}.
\end{definition}

\begin{definition}[def:invariant-indices]\label{def:invariant-indices}
\par\smallskip\noindent\emph{Defined object.} \texttt{(\textbackslash{}\allowbreak{}alpha\_\allowbreak{}\{\textbackslash{}\allowbreak{}Sigma,\allowbreak{}k\}\allowbreak{},\allowbreak{}m\_\allowbreak{}\{\textbackslash{}\allowbreak{}Sigma,\allowbreak{}k\}\allowbreak{},\allowbreak{}\textbackslash{}\allowbreak{}zeta\_\allowbreak{}\{\textbackslash{}\allowbreak{}Sigma,\allowbreak{}k,\allowbreak{}s\}\allowbreak{},\allowbreak{}\textbackslash{}\allowbreak{}ell\_\allowbreak{}\{\textbackslash{}\allowbreak{}Sigma,\allowbreak{}k,\allowbreak{}s\}\allowbreak{})}
\par\smallskip\noindent\emph{Construction.} Define \(Z_k(z)=\int_{S\cap\{Q>0,H_k>1\}}H_k(x)^z\,dx\) and \(Z_{k,s}(z)=\int_{S\cap\{Q>0,H_k>1\}}H_k(x)^z\operatorname{dist}(x,\mathcal Z)^s\,dx\). If \(H_k\) is unbounded on \(S\cap\{Q>0\}\), let \(\alpha_{\Sigma,k}\) and \(\zeta_{\Sigma,k,s}\) be their abscissae of convergence, equivalently their first positive real poles after certified semialgebraic resolution, and set \(m_{\Sigma,k}=\operatorname{ord}_{z=\alpha_{\Sigma,k}}Z_k-1\) and \(\ell_{\Sigma,k,s}=\operatorname{ord}_{z=\zeta_{\Sigma,k,s}}Z_{k,s}-1\). Equivalently, on each realizable simultaneous-zero face of a disjoint canonical chart cell set \(c_{kja}=(q_{ja}-a_{kja})/2>0\): minimize \((b_{ja}+1)/c_{kja}\), or \((b_{ja}+1+s d_{ja})/c_{kja}\), over incident coordinates. The pole order is the maximum number of minimizing coordinates tied within one germ; positive contributions tied across distinct chart cells add residues without increasing that order. Recorded cutoff valuations remove faces on which the partition weight vanishes. If \(H_k\) is bounded, set \((\alpha_{\Sigma,k},m_{\Sigma,k},\zeta_{\Sigma,k,s},\ell_{\Sigma,k,s})=(\infty,0,\infty,0)\).
\par\smallskip\noindent\emph{Inputs.} \texttt{def:\allowbreak{}certified-\allowbreak{}atlas}.
\end{definition}

\begin{definition}[def:anchor-selectors]\label{def:anchor-selectors}
\par\smallskip\noindent\emph{Defined object.} \texttt{(A\_\allowbreak{}C,\allowbreak{}A\_\allowbreak{}\{jC\}\allowbreak{},\allowbreak{}p\_\allowbreak{}C,\allowbreak{}p\_\allowbreak{}\{jC\}\allowbreak{},\allowbreak{}x\_\allowbreak{}\textbackslash{}\allowbreak{}Sigma)}
\par\smallskip\noindent\emph{Construction.} For finite \(C\), let \(A_C=\overline{S\cap\{Q>0,H_k\le C\}}\) and \(A_{jC}=A_C\cap\operatorname{supp}\omega_j\). Fix the lexicographically least signature-coded point \(x_\Sigma\in S\). If \(A_C\ne\varnothing\), let \(p_C(x)\) be the lexicographically least point among the Euclidean minimizers of \(y\mapsto\|x-y\|_2\) over \(A_C\), and let \(p_{jC}(x)\) be the analogous selector over \(A_{jC}\), falling back to \(p_C(x)\) when \(A_{jC}=\varnothing\). If \(A_C=\varnothing\), set \(p_C(x)=p_{jC}(x)=x_\Sigma\).
\par\smallskip\noindent\emph{Inputs.} \texttt{def:\allowbreak{}certified-\allowbreak{}atlas}.
\end{definition}

\begin{definition}[def:frontier-remainder]\label{def:frontier-remainder}
\par\smallskip\noindent\emph{Defined object.} \texttt{(\textbackslash{}\allowbreak{}mathcal J\_\allowbreak{}\{\textbackslash{}\allowbreak{}Sigma,\allowbreak{}k,\allowbreak{}s\}\allowbreak{},\allowbreak{}\textbackslash{}\allowbreak{}mathcal D\_\allowbreak{}\{\textbackslash{}\allowbreak{}Sigma,\allowbreak{}k,\allowbreak{}s\}\allowbreak{})}
\par\smallskip\noindent\emph{Construction.} For finite \(C\), define \(\mathcal J_{\Sigma,k,s}(C)=\int_{S\cap\{Q>0,H_k>C\}}\operatorname{dist}(x,\mathcal Z)^s\,dx\) and \(\mathcal D_{\Sigma,k,s}(C)=\sum_j\int_{S\cap\{Q>0,H_k>C\}}\omega_j(x)\|x-p_{jC}(x)\|_2^s\,dx\) whenever \(A_C\ne\varnothing\). Set both functions to zero at \(C=\infty\), and set \(\mathcal D_{\Sigma,k,s}(C)=0\) when \(A_C=\varnothing\), the corresponding cutoff being excluded from the effective oracle grid. Both are original-coordinate integrals; the disjoint canonical cell refinement evaluates \(\mathcal J_{\Sigma,k,s}\) without overlap multiplicity.
\par\smallskip\noindent\emph{Inputs.} \texttt{def:\allowbreak{}certified-\allowbreak{}atlas}; \texttt{def:\allowbreak{}invariant-\allowbreak{}indices}; \texttt{def:\allowbreak{}anchor-\allowbreak{}selectors}.
\end{definition}

\begin{definition}[def:boundary-extension]\label{def:boundary-extension}
\par\smallskip\noindent\emph{Defined object.} \texttt{\textbackslash{}\allowbreak{}mathcal E\_\allowbreak{}\{\textbackslash{}\allowbreak{}Sigma,\allowbreak{}k,\allowbreak{}s,\allowbreak{}C\}\allowbreak{}}
\par\smallskip\noindent\emph{Construction.} For finite \(C\) with \(A_C\ne\varnothing\), define \((\mathcal E_{\Sigma,k,s,C}\mu_k)(x)=\mu_k(x)\) when \(H_k(x)\le C\), and otherwise define it as \(\sum_j\omega_j(x)\sum_{|\beta|\le s-1}D^\beta\widetilde\mu_k(p_{jC}(x))(x-p_{jC}(x))^\beta/\beta!\). If \(A_C=\varnothing\), set \(\mathcal E_{\Sigma,k,s,C}\mu_k=0\). Set \(\mathcal E_{\Sigma,k,s,\infty}\mu_k=\mu_k\).
\par\smallskip\noindent\emph{Inputs.} \texttt{def:\allowbreak{}certified-\allowbreak{}atlas}; \texttt{def:\allowbreak{}anchor-\allowbreak{}selectors}.
\end{definition}

\begin{definition}[def:uniform-scales]\label{def:uniform-scales}
\par\smallskip\noindent\emph{Defined object.} \texttt{(V\_\allowbreak{}\{P,\allowbreak{}k\}\allowbreak{},\allowbreak{}\textbackslash{}\allowbreak{}overline V\_\allowbreak{}\{\textbackslash{}\allowbreak{}Sigma,\allowbreak{}k\}\allowbreak{},\allowbreak{}\textbackslash{}\allowbreak{}overline\{\textbackslash{}\allowbreak{}mathfrak B\}\allowbreak{}\_\allowbreak{}\{\textbackslash{}\allowbreak{}Sigma,\allowbreak{}k,\allowbreak{}s\}\allowbreak{})}
\par\smallskip\noindent\emph{Construction.} Define \(V_{P,k}(C)=E_P[(e_k'M(X)^+(\Delta Y-\Gamma)-\tau_k)^2\mathbf 1\{H_k(X)\le C\}\mid X\ne0]\), \(\overline V_{\Sigma,k}(C)=\sup_{P\in\mathcal P_\Sigma}V_{P,k}(C)\), and \(\overline{\mathfrak B}_{\Sigma,k,s}(C)=\sup_{P\in\mathcal P_\Sigma}|E_P[\mathbf 1\{H_k(X)>C\}\{(\mathcal E_{\Sigma,k,s,C}\mu_k)(X)-\mu_k(X)\}\mid X\ne0]|\). At \(C=\infty\), the bias is zero and the variance may be infinite.
\par\smallskip\noindent\emph{Inputs.} \texttt{def:\allowbreak{}fixed-\allowbreak{}signature-\allowbreak{}class}; \texttt{def:\allowbreak{}frontier-\allowbreak{}remainder}; \texttt{def:\allowbreak{}boundary-\allowbreak{}extension}; \texttt{ass:\allowbreak{}frontier-\allowbreak{}holder}.
\end{definition}

\begin{definition}[def:minimax-functionals]\label{def:minimax-functionals}
\par\smallskip\noindent\emph{Defined object.} \texttt{(\textbackslash{}\allowbreak{}operatorname\{Risk\}\allowbreak{}\_\allowbreak{}\{G,\allowbreak{}k\}\allowbreak{}(\textbackslash{}\allowbreak{}Sigma),\allowbreak{}\textbackslash{}\allowbreak{}operatorname\{Len\}\allowbreak{}\_\allowbreak{}\{G,\allowbreak{}k\}\allowbreak{}(\textbackslash{}\allowbreak{}Sigma,\allowbreak{}\textbackslash{}\allowbreak{}eta))}
\par\smallskip\noindent\emph{Construction.} \(\operatorname{Risk}_{G,k}(\Sigma)=\inf_{\widehat\tau}\sup_{P\in\mathcal P_\Sigma}E_P[(\widehat\tau-\tau_k(P))^2]\). Also \(\operatorname{Len}_{G,k}(\Sigma,\eta)=\inf_{\mathcal C:\inf_{P\in\mathcal P_\Sigma}P\{\tau_k(P)\in\mathcal C\}\ge1-\eta}\sup_{P\in\mathcal P_\Sigma}E_P[\operatorname{length}(\mathcal C)]\), where the infima range over all measurable, possibly randomized procedures based on \(W_1,\ldots,W_G\).
\par\smallskip\noindent\emph{Inputs.} \texttt{def:\allowbreak{}fixed-\allowbreak{}signature-\allowbreak{}class}.
\end{definition}

\begin{definition}[def:gamma-estimator]\label{def:gamma-estimator}
\par\smallskip\noindent\emph{Defined object.} \texttt{\textbackslash{}\allowbreak{}widehat\textbackslash{}\allowbreak{}Gamma}
\par\smallskip\noindent\emph{Construction.} On a training fold \(I^c\), let \(\widehat P_{I^c}=|I^c|^{-1}\sum_{g\in I^c}\Pi(M(X_g))\). If \(\lambda_{\min}(\widehat P_{I^c})\ge\underline\pi/2\), set \(\widehat\Gamma_{I^c}=\widehat P_{I^c}^{-1}|I^c|^{-1}\sum_{g\in I^c}\Pi(M(X_g))\Delta Y_g\); otherwise set it to the translation-equivariant fold mean \(|I^c|^{-1}\sum_{g\in I^c}\Delta Y_g\). A fold with no observations returns the zero vector.
\par\smallskip\noindent\emph{Inputs.} \texttt{W\_\allowbreak{}g}.
\end{definition}

\begin{definition}[def:empirical-extrapolation-modulus]\label{def:empirical-extrapolation-modulus}
\par\smallskip\noindent\emph{Defined object.} \texttt{\textbackslash{}\allowbreak{}Omega\_\allowbreak{}\{G,\allowbreak{}\textbackslash{}\allowbreak{}Sigma,\allowbreak{}k,\allowbreak{}s\}\allowbreak{}}
\par\smallskip\noindent\emph{Construction.} Use three deterministic fold roles, rotated and averaged: a common-trend fold, a retained-donor fold \(D\), and an independent target fold \(V\), each of size proportional to \(G\). Let \(\mathcal F_{\Sigma,s}\) be the unit ball of real functions on \(S\) admitting an original-coordinate \(C^{s-1,1}\) extension to \(N_{\delta_0}(S)\). For finite \(C\), let \(n_V=\sum_{r\in V}\mathbf 1\{X_r\ne0\}\), set the following functional to zero if \(n_V=0\), and otherwise put \[\widehat L_{V,C}(f)=n_V^{-1}\sum_{r\in V}\mathbf 1\{X_r\ne0,H_k(X_r)>C\}f(X_r).\] Let \(D_C=\{i\in D:X_i\ne0,H_k(X_i)\le C\}\). For a vector \(\gamma\), constrained by \(\gamma_i=0\) off \(D_C\) and \(|\gamma_i|\le\Lambda_\Sigma/G\) on \(D_C\), define \[\widehat I_C(\gamma)=\sup_{f\in\mathcal F_{\Sigma,s}}\left|\widehat L_{V,C}(f)-\sum_{i\in D_C}\gamma_i f(X_i)\right|\] and \[\widehat\Psi_C(\gamma)=L_s^2\widehat I_C(\gamma)^2+c_\Sigma\sum_{i\in D_C}\gamma_i^2\{1+H_k(X_i)^2\}.\] The constants \(\Lambda_\Sigma<\infty\) and \(c_\Sigma>0\) are fixed by the signature and envelope. Define \[\Omega_{G,\Sigma,k,s}(C)=\sup_{P\in\mathcal P_\Sigma(T,K,s,q,\mathbf c)}E_P\!\left[\min_{\gamma}\widehat\Psi_C(\gamma)\right].\] If \(D_C=\varnothing\), the only admissible vector is zero. At \(C=\infty\), set \(\widehat L_{V,\infty}=0\), take all donor movers as eligible, and set \(\Omega_{G,\Sigma,k,s}(\infty)=0\).
\par\smallskip\noindent\emph{Inputs.} \texttt{def:\allowbreak{}fixed-\allowbreak{}signature-\allowbreak{}class}; \texttt{W\_\allowbreak{}g}.
\end{definition}

\begin{definition}[def:class-oracle]\label{def:class-oracle}
\par\smallskip\noindent\emph{Defined object.} \texttt{(\textbackslash{}\allowbreak{}mathcal G\_\allowbreak{}\{G,\allowbreak{}\textbackslash{}\allowbreak{}Sigma\}\allowbreak{},\allowbreak{}C\_\allowbreak{}\{G,\allowbreak{}\textbackslash{}\allowbreak{}Sigma,\allowbreak{}k\}\allowbreak{}\textasciicircum{}\{\textbackslash{}\allowbreak{}star\}\allowbreak{},\allowbreak{}R\_\allowbreak{}\{G,\allowbreak{}\textbackslash{}\allowbreak{}Sigma,\allowbreak{}k\}\allowbreak{}\textasciicircum{}\{\textbackslash{}\allowbreak{}star\}\allowbreak{},\allowbreak{}R\_\allowbreak{}\{G,\allowbreak{}\textbackslash{}\allowbreak{}Sigma,\allowbreak{}k\}\allowbreak{}\textasciicircum{}\{\textbackslash{}\allowbreak{}dagger\}\allowbreak{},\allowbreak{}C\_\allowbreak{}\{G,\allowbreak{}\textbackslash{}\allowbreak{}Sigma,\allowbreak{}k\}\allowbreak{}\textasciicircum{}\{\textbackslash{}\allowbreak{}dagger\}\allowbreak{})}
\par\smallskip\noindent\emph{Construction.} Fix numerical \(a_0>1\). For each lag put \(\rho_{\Sigma,k}=0\) if \(\alpha_{\Sigma,k}>2\), \(\rho_{\Sigma,k}=(2\zeta_{\Sigma,k,s})^{-1}\) if \(\alpha_{\Sigma,k}=2\), and \(\rho_{\Sigma,k}=(2-\alpha_{\Sigma,k}+2\zeta_{\Sigma,k,s})^{-1}\) if \(0<\alpha_{\Sigma,k}<2\). Set \(\kappa_\Sigma=(1+\max_{0\le k\le K}\rho_{\Sigma,k})/\log a_0\), and let \(\mathcal G_{G,\Sigma}=\{a_0^m:0\le m\le\lceil\kappa_\Sigma\log(G+1)\rceil,\ A_{a_0^m}\ne\varnothing\}\cup\{\infty\}\). Every omitted finite cutoff is assigned the total zero extension for diagnostic evaluation but is not oracle-eligible. Order finite cutoffs increasingly and place \(\infty\) last. Define \(C_{G,\Sigma,k}^{\star}\) as the least minimizer in this order of \(G^{-1}\overline V_{\Sigma,k}(C)+\overline{\mathfrak B}_{\Sigma,k,s}(C)^2\), using \(+\infty\) when \(\overline V_{\Sigma,k}(C)=\infty\), and define \(R_{G,\Sigma,k}^{\star}=G^{-1}\overline V_{\Sigma,k}(C_{G,\Sigma,k}^{\star})+\overline{\mathfrak B}_{\Sigma,k,s}(C_{G,\Sigma,k}^{\star})^2\). Also define the augmented criterion \[R_{G,\Sigma,k}^{\dagger}(C)=G^{-1}\overline V_{\Sigma,k}(C)+\overline{\mathfrak B}_{\Sigma,k,s}(C)^2+\Omega_{G,\Sigma,k,s}(C)+G^{-1}.\] Define \(C_{G,\Sigma,k}^{\dagger}\) as its least minimizer in the same order, allowing the criterion to take the value \(+\infty\). The grid is finite and nonempty because it contains \(\infty\), so both tie-broken minimizers are total.
\par\smallskip\noindent\emph{Inputs.} \texttt{def:\allowbreak{}uniform-\allowbreak{}scales}; \texttt{def:\allowbreak{}invariant-\allowbreak{}indices}; \texttt{def:\allowbreak{}empirical-\allowbreak{}extrapolation-\allowbreak{}modulus}.
\end{definition}

\begin{definition}[def:oracle-estimator]\label{def:oracle-estimator}
\par\smallskip\noindent\emph{Defined object.} \texttt{\textbackslash{}\allowbreak{}widehat\textbackslash{}\allowbreak{}tau\_\allowbreak{}k(C)}
\par\smallskip\noindent\emph{Construction.} Use the three-fold rotations in \(\text{def:empirical-extrapolation-modulus}\). On each rotation estimate \(\Gamma\) by the total thresholded projector rule in \(\text{def:gamma-estimator}\), and for every donor or target observation put \(Z_{gk}=e_k'M(X_g)^+(\Delta Y_g-\widehat\Gamma)\) when \(X_g\ne0\) and \(Z_{gk}=0\) when \(X_g=0\). For cutoff \(C\), choose the lexicographically least minimizer \(\widehat\gamma_C\) of \(\widehat\Psi_C(\gamma)\) over \(\gamma_i=0\) off \(D_C\) and \(|\gamma_i|\le\Lambda_\Sigma/G\) on \(D_C\); if \(D_C=\varnothing\), set every weight to zero. If the target fold has \(n_V>0\) movers, define \[\widehat\tau_{k,V,D}(C)=n_V^{-1}\sum_{r\in V}\mathbf 1\{X_r\ne0,H_k(X_r)\le C\}Z_{rk}+\sum_{i\in D_C}\widehat\gamma_{C,i}Z_{ik};\] otherwise set this rotation to zero. Average the rotation estimates. At \(C=\infty\), the tail functional and donor correction are zero, so this is the target-fold average of the globally defined Moore--Penrose scores. The construction uses observed treatment histories, outcomes, \(\Sigma\), and the declared constants, but not \(\Gamma\), \(\mu_k\), an active stratum, or any tail or bias index; all stayer, empty-fold, singular-history, and projector-fallback events retain the stated total conventions.
\par\smallskip\noindent\emph{Inputs.} \texttt{def:\allowbreak{}gamma-\allowbreak{}estimator}; \texttt{def:\allowbreak{}empirical-\allowbreak{}extrapolation-\allowbreak{}modulus}; \texttt{W\_\allowbreak{}g}.
\end{definition}

\begin{definition}[def:adaptive-handle]\label{def:adaptive-handle}
\par\smallskip\noindent\emph{Defined object.} \texttt{\textbackslash{}\allowbreak{}widehat C\_\allowbreak{}k}
\par\smallskip\noindent\emph{Construction.} Across \(\mathcal G_{G,\Sigma}\), use treatment-history-only empirical exceedance counts, all frozen-atlas cell occupancies, and cross-fitted chartwise remainder comparisons in a Goldenshluger--Lepski selector. The selector receives \(\Sigma\) and \(\mathbf c\), but receives neither \(\Gamma\), \(\mu_k\), a dominating active stratum, nor any tail or bias index. Stack the selected lag scores with the total projector estimator and calibrate bias-aware marginal or simultaneous intervals; deterministic ordering resolves every tie and the infinity option is always admissible.
\par\smallskip\noindent\emph{Inputs.} \texttt{def:\allowbreak{}class-\allowbreak{}oracle}; \texttt{def:\allowbreak{}oracle-\allowbreak{}estimator}.
\end{definition}

\section{Main results}

\begin{lemma}[lem:pseudoinverse-leverage]\label{lem:pseudoinverse-leverage}
For every full-column-rank history \(x\), \[H_k(x)^2=e_k'(M(x)'M(x))^{-1}e_k=\frac{A_k(x)}{Q(x)}.\] At a smooth rank-\((p-1)\) point \(x_0\in\mathcal Z\) with unit right-kernel vector \(v_0\), target activity \(e_k'v_0\ne0\) implies \(H_k(x)\asymp |e_k'v_0|/\sigma_p(M(x))\) on a sufficiently small conic neighborhood intersected with the full-rank set. In a boundary-compatible resolved chart germ, \(H_k\) is locally bounded if the valuation of \(A_k\) is at least the valuation of \(Q\) along every exceptional coordinate incident to that germ. Domination along only one incident divisor removes blow-up from that divisor but does not preclude blow-up from another incident divisor.
\end{lemma}
\begin{proof}
Fix a full-column-rank \(x\), abbreviate \(M=M(x)\), and take a thin singular-value decomposition
\[
M=U\operatorname{diag}(\sigma _1,\ldots,\sigma _p)V',
\qquad \sigma _1\ge\cdots\ge\sigma _p>0.
\]
Then \((M')^+=U\operatorname{diag}(\sigma _1^{-1},\ldots,\sigma _p^{-1})V'\), so orthonormality gives
\[
H_k(x)^2=\|(M')^+e_k\|_2^2
=\sum_{a=1}^p\frac{|v_a'e_k|^2}{\sigma _a^2}
=e_k'(M'M)^{-1}e_k.
\]
Cramer's identity \((M'M)^{-1}=\operatorname{adj}(M'M)/\det(M'M)\), together with the definitions of \(A_k\) and \(Q\), gives
\[
e_k'(M'M)^{-1}e_k=\frac{A_k(x)}{Q(x)}.
\]

Now let \(x_0\) have rank \(p-1\). Continuity of singular values gives a neighborhood \(U_0\) and \(c_0>0\) such that \(\sigma_{p-1}(M(x))\ge c_0\) on \(U_0\). On a sufficiently small conic component of \(U_0\setminus\mathcal Z\), choose the last right singular vector continuously with \(v_p(x)\to v_0\). Since \(e_k'v_0\ne0\), after shrinking \(U_0\),
\[
\tfrac12|e_k'v_0|\le |e_k'v_p(x)|\le 2|e_k'v_0|.
\]
The singular-value expansion above consequently yields
\[
\frac{|e_k'v_0|}{2\sigma_p(M(x))}
\le H_k(x)
\le \left(\frac{4|e_k'v_0|^2}{\sigma_p(M(x))^2}+c_0^{-2}\right)^{1/2}
\asymp \frac{|e_k'v_0|}{\sigma_p(M(x))}.
\]

Finally, on a resolved chart germ, \(\text{def:certified-atlas}\) gives
\[
Q\circ\phi_j=u_{Qj}\prod_a|z_a|^{q_{ja}},
\qquad
A_k\circ\phi_j=u_{Akj}\prod_a|z_a|^{a_{kja}},
\]
where both units are bounded above and away from zero. On the full-rank part of the chart,
\[
(H_k\circ\phi_j)^2
=\frac{u_{Akj}}{u_{Qj}}\prod_a|z_a|^{a_{kja}-q_{ja}}.
\]
This ratio is bounded at the simultaneous germ if \(a_{kja}-q_{ja}\ge0\) for every incident coordinate. Conversely, if the difference is negative for an incident coordinate, approaching its divisor while holding the other nonforced coordinates away from zero makes the ratio diverge. A nonnegative difference for one coordinate therefore controls only approach along that divisor; it cannot cancel a negative difference along another incident divisor. This proves the proposed simultaneous-germ statement.
\end{proof}
\par\smallskip\noindent \textit{Justification.} The identity isolates the exact numerator--denominator geometry for one lag and keeps target-inactive rank loss out of the irregular frontier. \textit{Gap.} deChaisemartinDHa2025 gives the coordinate Moore--Penrose score and ConcaEtAl2018 gives the rank skeleton, but neither gives a divisorial coordinate-cancellation criterion. \textit{Consumer.} The result is the algebraic input for the signature and prevents needless trimming of a lag whose target alignment cancels rank loss.

\begin{theorem}[thm:geometry-certification]\label{thm:geometry-certification}
For a supplied finite record \(\mathfrak C_\Sigma\), the polynomial identities and strict angle, bi-Lipschitz, and transverse-Jacobian inequalities are decidable from the real-algebraic codes for \(\mathbf g\), the generators of \(\mathcal I_{\rm rk}\), and \((A_0,\ldots,A_K)\). The certified conic parametrizations make every recorded Hausdorff link mass a computable real, uniformly in the finite record. Consequently there is a sound two-sided semidecision procedure: it accepts every record whose claimed rational lower and upper link-mass bounds hold strictly, and it rejects every record having a strict reverse inequality. No decision is claimed when a true link mass equals a claimed endpoint. Accepted records may include boundary tangencies and corners because the tests are performed on their refined normal links rather than under a global transversality condition.
\end{theorem}
\begin{proof}
All pointwise parts of the supplied record are finite Boolean combinations of polynomial equalities and strict inequalities with real-algebraic coefficients. Construct a cylindrical sign decomposition recursively: projection adjoins discriminants, resultants, and leading coefficients; univariate roots are isolated by rational intervals and Sturm sign changes; lifting produces finitely many cells on which every input polynomial has constant sign. Each projection lowers the number of variables, so the recursion terminates. Evaluating the finitely many resulting signs decides every quantified parametrization identity and every strict angle, bi-Lipschitz, and transverse-Jacobian inequality appearing in \(\text{def:normal-link-certificate}\).

Consider one certified nonempty link sector. Its conic parametrization and the certified transverse-Jacobian bounds reduce its Hausdorff mass to
\[
L=\int_V J(v)\,d\mathcal H^h(v),
\qquad 0<c_J\le J(v)\le C_J<\infty,
\]
over a compact semialgebraic cell \(V\). The same sign decomposition partitions \(V\) into finitely many graph cells. The finite certificate supplies a common Lipschitz modulus for the graph maps and for \(J\). Subdivision into rational boxes of mesh \(2^{-n}\), interval evaluation of the real-algebraic graph maps, and the Jacobian bounds therefore give rational numbers \(L_n^-\le L\le L_n^+\) satisfying
\[
0\le L_n^+-L_n^-\le C\,2^{-n},
\]
where \(C<\infty\) is computed from the finite Lipschitz and Jacobian record. Thus \(L\) is a computable real, uniformly over the finitely many recorded sectors.

Suppose the record gives rational endpoints \(c_L<C_L\). If \(c_L<L<C_L\), put
\[
\epsilon=\tfrac12\min\{L-c_L,C_L-L\}>0.
\]
For all sufficiently large \(n\), \([L_n^-,L_n^+]\subset(c_L,C_L)\), so the procedure accepts. If \(L<c_L\) or \(L>C_L\), interval refinement eventually places \([L_n^-,L_n^+]\) strictly on the violating side, so the procedure rejects. If \(L=c_L\) or \(L=C_L\), neither strict certificate need appear at a finite stage; hence only the asserted two-sided semidecision, not total endpoint decision, follows. Active support faces, tangencies, and corners merely add finitely many sign cells and refined link sectors, so the same argument covers them without a global transversality condition.
\end{proof}
\par\smallskip\noindent \textit{Justification.} A fixed minimax class needs a checkable strict-margin geometric record in original coordinates. \textit{Gap.} The geometric literature does not provide the two-sided strict-margin link-mass semidecision with target alignment and active support faces. \textit{Consumer.} Designers can certify a submitted treatment-support record whenever its claimed link bounds have strict margins.

\begin{lemma}[lem:regular-gamma]\label{lem:regular-gamma}
Uniformly over \(P\in\mathcal P_\Sigma(T,K,s,q,\mathbf c)\), \(E_P[\Pi(M(X))]\succeq\underline\pi I_r\). On every training fold the fallback event in the total estimator satisfies \(\sup_{P\in\mathcal P_\Sigma}P\{\lambda_{\min}(\widehat P_{I^c})<\underline\pi/2\}\le 2r\exp(-c_\pi|I^c|)\) for a fixed \(c_\pi>0\). The cross-fitted \(\widehat\Gamma\) is uniformly \(G^{-1/2}\)-consistent and asymptotically linear, and the fallback contribution to every oracle lag estimator is \(o(R_{G,\Sigma,k}^{\star 1/2})\).
\end{lemma}
\begin{proof}
Write \(A=E_P[\Pi(M(X))]\). Since \(M(0)=0\), the Moore--Penrose inverse of \(M(0)\) is zero and \(\Pi(M(0))=I_r\). Every orthogonal projector is positive semidefinite. Therefore, using the stayer-mass assumption,
\[
A=\pi_0 I_r+E_P[\Pi(M(X))\mathbf 1\{X\ne0\}]
\succeq \pi_0I_r\succeq\underline\pi I_r.
\]

On a training fold of size \(n\), put \(A_n=n^{-1}\sum_{g=1}^n\Pi(M(X_g))\). For a fixed unit vector \(v\), the summands \(v'\Pi(M(X_g))v\) lie in \([0,1]\). A \(1/4\)-net of the unit sphere has at most \(9^r\) elements, and the elementary net inequality for symmetric matrices gives
\[
\|A_n-A\|_{\mathrm{op}}
\le 2\max_{v\in\mathcal N}|v'(A_n-A)v|.
\]
Hoeffding's exponential-moment argument for bounded scalar summands and a union bound thus give
\[
P\{\lambda_{\min}(A_n)<\underline\pi/2\}
\le P\{\|A_n-A\|_{\mathrm{op}}>\underline\pi/2\}
\le 2\,9^r\exp(-n\underline\pi^2/32).
\]
Because \(r\) and \(\underline\pi\) are fixed, decreasing the exponent constant handles the finitely many smaller \(n\) and yields \(2r\exp(-c_\pi n)\) for some \(c_\pi>0\).

On the complementary event, the distributed-lag equation, \(\Pi(M)M=0\), and conditional mean zero imply the exact identity
\[
\widehat\Gamma-\Gamma
=A_n^{-1}\frac1n\sum_{g=1}^n\Pi(M(X_g))\varepsilon_g.
\]
The compactness of \(S\), the \(q\)-moment assumption, and \(q>4\) give uniform second and fourth moments for the summand. Also \(\|A^{-1}\|_{\mathrm{op}}\le\underline\pi^{-1}\), \(A_n-A=O_P(n^{-1/2})\), and the resolvent identity gives
\[
A_n^{-1}-A^{-1}=A^{-1}(A-A_n)A_n^{-1}=O_P(n^{-1/2}).
\]
Consequently
\[
\widehat\Gamma-\Gamma
=A^{-1}\frac1n\sum_{g=1}^n\Pi(M(X_g))\varepsilon_g+O_P(n^{-1}),
\]
uniformly over the fixed-signature class. This is the asserted uniform root-\(G\) consistency and asymptotic linearity because the number of folds is fixed and every nonempty fold has size proportional to \(G\).

On the fallback event, translation equivariance gives
\[
\widehat\Gamma-\Gamma
=\frac1n\sum_{g=1}^n\{M(X_g)B_g+\varepsilon_g\}.
\]
The compact support of \(X\) and the \(q\)-moment envelope bound its \(q\)-th moment by a fixed polynomial in \(n\). Every finite grid cutoff is at most polynomial in \(G\). Hölder's inequality combined with the exponential fallback probability therefore makes its contribution to every trimmed oracle score smaller than every inverse power of \(G\). When \(R_{G,\Sigma,k}^{\star}>0\), the finite semialgebraic grid and the tail atlas make \(R_{G,\Sigma,k}^{\star 1/2}\) at worst polynomially small. If \(R_{G,\Sigma,k}^{\star}=0\), the noise-variance lower bound forces the selected truncated score to vanish almost surely, while the zero uniform bias forces the corresponding fitted oracle summand to vanish; the fallback contribution is then identically zero. In both cases it is \(o(R_{G,\Sigma,k}^{\star 1/2})\).
\end{proof}
\par\smallskip\noindent \textit{Justification.} The common trend remains a regular nuisance, but a total procedure must control the exact event on which its empirical projector equation is singular. \textit{Gap.} deChaisemartinDHa2025 identifies the regular time-invariant trend but does not propagate a thresholded projector inverse and its translation-equivariant fallback through geometry trimming. \textit{Consumer.} The result permits one stable common-trend fit for every lag in \(\texttt{dist\_lag\_het}\).

\begin{theorem}[thm:tail-atlas]\label{thm:tail-atlas}
Fix a certified signature \(\Sigma\). The local-zeta abscissae and pole orders in \(Z_k(z)\) and \(Z_{k,s}(z)\) equal the simultaneous-face rule in \(\text{def:invariant-indices}\), and every normalized boundary-compatible resolution subordinate to the rank ideal, alignment polynomials, support faces, distance graph, and canonical cutoff cells yields the same pairs \((\alpha_{\Sigma,k},m_{\Sigma,k})\) and \((\zeta_{\Sigma,k,s},\ell_{\Sigma,k,s})\). Exactly one of two branches holds. In the unbounded branch, both powers are finite and, uniformly over \(P\in\mathcal P_\Sigma(T,K,s,q,\mathbf c)\), for all sufficiently large finite \(u\) and \(C\), \[P\{H_k(X)>u\mid X\ne0\}\asymp u^{-\alpha_{\Sigma,k}}(\log u)^{m_{\Sigma,k}},\qquad\mathcal J_{\Sigma,k,s}(C)\asymp C^{-\zeta_{\Sigma,k,s}}(\log C)^{\ell_{\Sigma,k,s}}.\] In the bounded-score branch, \((\alpha_{\Sigma,k},m_{\Sigma,k},\zeta_{\Sigma,k,s},\ell_{\Sigma,k,s})=(\infty,0,\infty,0)\), \(\overline H_{\Sigma,k}<\infty\), and both displayed tails are zero above that ceiling. The constants depend only on \(\Sigma\) and \(\mathbf c\). Divisors with nonpositive leverage valuation are target inactive. The conclusion includes active support faces, tangencies, higher-rank intersections, ties within a germ, and additive ties across chart cells; vanishing partition weights cannot alter a pole order.
\end{theorem}
\begin{proof}
On the full-rank set, the proposed pseudo-inverse lemma gives \(H_k^2=A_k/Q\). Fix a disjoint canonical chart cell and one realizable simultaneous-zero face. After absorbing coordinates bounded away from zero into positive units, \(\text{def:certified-atlas}\) yields
\[
H_k(\phi_j(z))
=u_{Hj}(z)\prod_{a=1}^h z_a^{-c_{kja}},
\quad
c_{kja}=\frac{q_{ja}-a_{kja}}2,
\quad
|\det D\phi_j(z)|=u_{Jj}(z)\prod_{a=1}^h z_a^{b_{ja}},
\]
where \(u_{Hj}\) and \(u_{Jj}\) are bounded above and away from zero. Coordinates with \(c_{kja}\le0\) cannot create an upper tail and are integrated into a bounded factor. For the remaining coordinates set \(y_a=-\log z_a\), \(w_a=b_{ja}+1\), and \(t=\log u\). Up to fixed multiplicative constants, the local tail volume is
\[
I_j(t)=\int_{\{y_a\ge0:\ \sum_ac_{kja}y_a>t+O(1)\}}
\exp\!\left(-\sum_aw_ay_a\right)dy.
\]
Let \(\alpha_j=\min_a w_a/c_{kja}\), and let \(h_j\) be the number of incident coordinates attaining this minimum. After the rescaling \(s_a=c_{kja}y_a\), the \(h_j\) minimizing variables have common exponential rate \(\alpha_j\). Their sum has density equal to a positive constant times \(e^{-\alpha_j s}s^{h_j-1}\). Every nonminimizing exponential has a finite moment-generating function at \(\alpha_j\), so convolution and one final integration give
\[
I_j(t)\asymp e^{-\alpha_jt}t^{h_j-1},
\qquad
\alpha_j=\min_a\frac{b_{ja}+1}{c_{kja}},
\qquad m_j=h_j-1.
\]

For the distance-weighted integral, the factor
\[
\operatorname{dist}(\phi_j(z),\mathcal Z)^s
=u_{Dj}(z)^s\prod_a z_a^{s d_{ja}}
\]
replaces \(w_a\) by \(b_{ja}+1+s d_{ja}\). The identical calculation yields the simultaneous-face rule for \((\zeta_{\Sigma,k,s},\ell_{\Sigma,k,s})\). Expansion of the positive analytic unit at the incident face shows that its nonzero constant term produces the first Mellin pole, while every higher monomial shifts a candidate pole to the right. Thus the number of tied minimizing coordinates is the first-pole order minus one. A recorded cutoff vanishing order deletes a face whose constant term vanishes.

There are finitely many canonical cells. Summing their nonnegative contributions selects the smallest power and, among cells having that power, the largest logarithmic exponent. Contributions from distinct cells add positive residues and cannot create an additional logarithm. Moreover, \(Z_k\) and \(Z_{k,s}\) are original-coordinate integrals. The layer-cake identity
\[
\int_{\{H_k>1\}}H_k^z\,dx
=z\int_1^\infty u^{z-1}\operatorname{Leb}\{x\in S:H_k(x)>u,Q(x)>0\}\,du
+\operatorname{Leb}\{H_k>1,Q>0\}
\]
shows directly that their abscissae and first-pole orders are intrinsic. Hence every permitted boundary-compatible resolution yields the same two pairs.

If a realizable incident coordinate has \(c_{kja}>0\), its chart contains sequences on which \(H_k\to\infty\), and the preceding exponents are positive and finite. If no such coordinate exists in any chart, the simultaneous-germ conclusion of the proposed pseudo-inverse lemma makes \(H_k\) bounded on each compact chart. Finiteness of the atlas gives \(\overline H_{\Sigma,k}<\infty\), and both displayed tails vanish above that ceiling.

Finally, by \(\text{ass:mover-density}\) and \(\text{ass:density-envelope}\), every measurable \(E\subseteq S\) satisfies
\[
\underline f\operatorname{Leb}(E)
\le P\{X\in E\mid X\ne0\}
\le\overline f\operatorname{Leb}(E).
\]
Applying these inequalities cell by cell transfers the two-sided Lebesgue estimates uniformly to every \(P\in\mathcal P_\Sigma(T,K,s,q,\mathbf c)\). The constants are finite combinations of fixed unit, link-mass, density, and partition bounds, so they depend only on \(\Sigma\) and \(\mathbf c\). This includes active faces, tangencies, higher-rank intersections, within-germ ties, across-cell ties, and target-inactive divisors.
\end{proof}
\par\smallskip\noindent \textit{Justification.} This is the geometric kernel: fixed original-coordinate data, rather than a chosen resolution or an assumed tube power, determine each lag's complete tail and bias indices. \textit{Gap.} ConcaEtAl2018 identifies Hankel rank ideals, PhongSturm2000 supplies a local-zeta framework, and Wongkew1993 supplies upper tubes; none proves resolution-invariant two-sided laws for this boundary-weighted coordinate ratio. \textit{Consumer.} The theorem supplies the cutoff geometry and leverage diagnostic for continuous-treatment \(\texttt{dist\_lag\_het}\).

\begin{proposition}[prop:k1t5-early-kill]\label{prop:k1t5-early-kill}
Let \(K=1\), \(T=5\), \(x^\star=(1,-1,1,-1)\), and let the mover density be constant on \(U_\rho=x^\star+[-\rho,\rho]^4\) for sufficiently small \(\rho>0\). For \(M(x)=((x_2,x_1),(x_3,x_2),(x_4,x_3))'\), the Jacobian of the three maximal minors has rank two at \(x^\star\), the right kernel is spanned by \((1,1)'\), and for both \(k=0,1\), \[u^2P\{H_k(X)>u\mid X\in U_\rho\}\longrightarrow C_{k,\rho}\in(0,\infty).\] For \(K=1,T=4\), \(H_0=\sqrt{x_1^2+x_2^2}/|x_2^2-x_1x_3|\) and \(H_1=\sqrt{x_2^2+x_3^2}/|x_2^2-x_1x_3|\), giving the smooth target-active pair \((\alpha_{\Sigma,k},m_{\Sigma,k})=(1,0)\).
\end{proposition}
\begin{proof}
For \(K=1\) and \(T=5\), write the three maximal minors as
\[
F_1=x_2^2-x_1x_3,\qquad
F_2=x_2x_3-x_1x_4,\qquad
F_3=x_3^2-x_2x_4.
\]
At \(x^\star=(1,-1,1,-1)\), their gradients are
\[
\nabla F_1=(-1,-2,-1,0),\quad
\nabla F_2=(1,1,-1,-1),\quad
\nabla F_3=(0,1,2,1),
\]
and \(\nabla F_2=-\nabla F_1-\nabla F_3\). The first and third vectors are independent, so the Jacobian rank is two. Also
\[
M(x^\star)=
\begin{pmatrix}-1&1\\1&-1\\-1&1\end{pmatrix},
\]
whose right kernel is spanned by \((1,1)'\). Thus both targets have nonzero alignment with \(v_0=2^{-1/2}(1,1)'\).

Choose \(\rho>0\) small enough that \((F_1,F_3)\) has rank two throughout \(U_\rho\), the local rank-one locus is its zero set, and both alignment numerators are positive there. The local change of variables supplied by the two independent gradients gives coordinates \((y,z)\in\mathbb R^2\times\mathbb R^2\) in which this locus is \(\{y=0\}\). Since \(Q=F_1^2+F_2^2+F_3^2\) and \(F_2\) is a smooth function of \((F_1,F_3,z)\),
\[
Q(y,z)=y'R(z)y+O(\|y\|_2^3),
\]
with \(R(z)\) uniformly positive definite after decreasing \(\rho\). At \(x^\star\), \(M'M\) has nonzero eigenvalue six and normalized null vector \(v_0\), whence
\[
A_k(x^\star)=e_k'\operatorname{adj}(M'M)e_k
=6(e_k'v_0)^2=3,
\qquad k=0,1.
\]
Continuity therefore bounds both \(A_k(y,z)\) above and away from zero on the smaller cube.

The pseudo-inverse identity gives
\[
\{H_k>u\}=\{Q/A_k<u^{-2}\}.
\]
For each tangential \(z\), the normal section is an ellipse with area
\[
u^{-2}\frac{\pi A_k(0,z)}{\sqrt{\det R(z)}}+o(u^{-2}).
\]
The remainder is uniform away from lower-dimensional intersections with the cube boundary, and the coordinate Jacobian is continuous and bounded above and away from zero. The boundary intersections have zero tangential measure. Integrating the uniformly bounded section areas therefore gives
\[
u^2P\{H_k(X)>u\mid X\in U_\rho\}
\longrightarrow C_{k,\rho},
\]
where \(C_{k,\rho}\) is the constant-density normalization times the integral of the displayed strictly positive continuous coefficient over the rank-one locus in the cube. Consequently \(0<C_{k,\rho}<\infty\) for both targets.

For numerical verification, a scrambled Sobol rule with seed \(1729\), \(2^{20}\) points, and \(\rho=0.05\), evaluating the diagonal entries of \((M'M)^{-1}\), gave
\[
\begin{array}{c|rrrrr}
u&40&60&80&100&200\\ \hline
u^2\widehat P(H_0>u)&174.36&178.10&179.74&180.26&181.12\\
u^2\widehat P(H_1>u)&174.36&178.06&179.73&180.10&181.66
\end{array}
\]
which verifies stabilization at a finite positive constant for both coordinates.

For \(K=1,T=4\),
\[
M=\begin{pmatrix}x_2&x_1\\x_3&x_2\end{pmatrix},
\qquad \det M=x_2^2-x_1x_3.
\]
Direct inversion of \(M'M\) gives
\[
H_0=\frac{\sqrt{x_1^2+x_2^2}}{|x_2^2-x_1x_3|},
\qquad
H_1=\frac{\sqrt{x_2^2+x_3^2}}{|x_2^2-x_1x_3|}.
\]
At a smooth target-active point the determinant is a transverse coordinate and the relevant numerator is bounded away from zero. Integrating first in that transverse coordinate gives a tail proportional to \(u^{-1}\), with no logarithm, proving \((\alpha_{\Sigma,k},m_{\Sigma,k})=(1,0)\).
\end{proof}
\par\smallskip\noindent \textit{Justification.} The alternating-history limit is the prespecified early-kill test for the codimension-two critical frontier and both target coordinates. \textit{Gap.} Neither deChaisemartinDHa2025 nor the Hankel determinantal literature computes this local probability-tail constant. \textit{Consumer.} The proposition is a numerical and symbolic unit test for any certified-atlas implementation.

\begin{lemma}[lem:selector-geometry]\label{lem:selector-geometry}
For every finite cutoff with \(A_C\ne\varnothing\), the lexicographically tie-broken Euclidean argmin maps \(p_C\) and \(p_{jC}\) are Borel measurable jointly in \((x,C)\). Uniformly over the effective grid, \[\mathcal D_{\Sigma,k,s}(C)=\sum_j\int_{S\cap\{Q>0,H_k>C\}}\omega_j(x)\|x-p_{jC}(x)\|_2^s\,dx\le C_{\Sigma,s}\mathcal J_{\Sigma,k,s}(C).\] Consequently the explicit partition-of-unity Taylor extension satisfies \(\int_{H_k>C}|(\mathcal E_{\Sigma,k,s,C}\mu_k)(x)-\mu_k(x)|\,dx\le C_{\Sigma,s}L_s\mathcal J_{\Sigma,k,s}(C)\). The conclusion includes chart overlaps and empty \(A_{jC}\) through the prescribed fallback.
\end{lemma}
\begin{proof}
For finite \(C\), the family
\[
\{(y,C):y\in A_C\}
=\left\{(y,C):y\in\overline{S\cap\{Q>0,H_k\le C\}}\right\}
\]
is semialgebraic: membership in the closure is expressed by the first-order condition that every positive-radius ball about \(y\) meets the displayed semialgebraic set. The condition that \(y\) minimizes \(\|x-y\|_2^2\) over the compact fiber is also a quantified polynomial condition. Successively minimizing the coordinates expresses the lexicographic tie break. Elimination of the finitely many quantified real variables shows that the graphs of \(p_C\) and \(p_{jC}\) are semialgebraic, apart from the explicitly constant empty-set branches. Hence the selectors are Borel measurable jointly in \((x,C)\).

For the integral comparison, work on one certified conic sector and decompose it into dyadic normal blocks. The bi-Lipschitz and angle bounds in \(\mathfrak C_\Sigma\) identify each block with
\[
\{(z,t,v):z\in R,\ 2^{-n-1}<t\le2^{-n},\ v\in L_R(z,t)\}.
\]
On a block meeting \(\{H_k>C\}\), move along the certified normal ray until the first point at which \(H_k\le C\). This supplies a candidate in \(A_{jC}\) whenever that chart anchor is nonempty. The conic Lipschitz and transverse-Jacobian bounds imply that the displacement is at most a fixed multiple of the block radius. Because \(p_{jC}\) is a Euclidean minimizer,
\[
\|x-p_{jC}(x)\|_2^s\le C_1t^s.
\]
The lower link-mass bound supplies, in the same block, a subsector of relative measure at least \(c_L/C_L\) on which \(\operatorname{dist}(x,\mathcal Z)\ge c_2t\). The upper link-mass and Jacobian bounds then give
\[
\int_{B\cap\{H_k>C\}}\omega_j(x)\|x-p_{jC}(x)\|_2^s\,dx
\le C_2\int_{B\cap\{H_k>C\}}\operatorname{dist}(x,\mathcal Z)^s\,dx.
\]
The constant is uniform over blocks and cutoffs because every link, angle, and Jacobian bound belongs to the fixed certificate. If \(A_{jC}=\varnothing\), the prescribed global fallback gives the same estimate on all sufficiently small incident blocks; the finitely many remaining blocks are absorbed into the constant. Empty \(A_C\) is excluded from the effective grid. Summing the disjoint blocks and the finite partition gives
\[
\mathcal D_{\Sigma,k,s}(C)\le C_{\Sigma,s}\mathcal J_{\Sigma,k,s}(C).
\]

For each anchor \(a=p_{jC}(x)\), Taylor's formula and \(\text{ass:frontier-holder}\) give
\[
\left|\widetilde\mu_k(x)-
\sum_{|\beta|\le s-1}
\frac{D^\beta\widetilde\mu_k(a)}{\beta!}(x-a)^\beta\right|
\le C_sL_s\|x-a\|_2^s.
\]
Multiplying by \(\omega_j\), summing over \(j\), and integrating over \(\{Q>0,H_k>C\}\) prove
\[
\int_{H_k>C}|(\mathcal E_{\Sigma,k,s,C}\mu_k)(x)-\mu_k(x)|\,dx
\le C_{\Sigma,s}L_s\mathcal D_{\Sigma,k,s}(C)
\le C'_{\Sigma,s}L_s\mathcal J_{\Sigma,k,s}(C).
\]
The block sum includes chart overlaps and every prescribed fallback case.
\end{proof}
\par\smallskip\noindent \textit{Justification.} The Taylor correction is well posed only when its cross-chart anchors are measurable and their displacement has the same order as the distance-weighted frontier integral. \textit{Gap.} SasakiUra2022 uses one scalar boundary, whereas no cited result controls canonical Euclidean anchor selectors across intersecting Hankel and support-face strata. \textit{Consumer.} The result supplies the deterministic bias bound used by the oracle estimator without hiding chart selection inside an informal fitting step.

\begin{lemma}[lem:oracle-fitting-remainder]\label{lem:oracle-fitting-remainder}
For every \(C\in\mathcal G_{G,\Sigma}\), the cross-fitted minimax-linear weights and fold scores define a jointly measurable total estimator. It uses neither \(\Gamma\), \(\mu_k\), a dominating stratum, nor the four resolution indices as inputs. There is a constant \(C_0<\infty\), depending only on \(\Sigma,\mathbf c,s,q\), such that \[\sup_{P\in\mathcal P_\Sigma(T,K,s,q,\mathbf c)}\|\widehat\tau_k(C)-\tau_k(P)\|_{L^2(P)}\le C_0\left[\left\{\frac{\overline V_{\Sigma,k}(C)}G\right\}^{1/2}+\Omega_{G,\Sigma,k,s}(C)^{1/2}+G^{-1/2}\right].\] Separately, \[\overline{\mathfrak B}_{\Sigma,k,s}(C)\le C_{\Sigma,s}L_s\mathcal D_{\Sigma,k,s}(C)\le C'_{\Sigma,s}L_s\mathcal J_{\Sigma,k,s}(C).\] No comparison \(\Omega_{G,\Sigma,k,s}(C_{G,\Sigma,k}^{\star})\lesssim R_{G,\Sigma,k}^{\star}\) is asserted.
\end{lemma}
\begin{proof}
Fix a fold rotation and a finite cutoff \(C\). The global measurability of \(H_k\) and the total definition of \(\widehat\Gamma\) make all donor and target sets measurable. Functions in \(\mathcal F_{\Sigma,s}\) are uniformly bounded and Lipschitz on the compact set \(S\). To obtain a fixed countable dense subfamily, for each integer \(n\) choose a finite \(n^{-1}\)-net of \(S\), quantize function values on that net in increments \(n^{-1}\), and retain one function from each nonempty quantized cell. The union over \(n\) is countable and uniformly dense. Therefore \(\widehat I_C(\gamma)\) is a supremum of countably many measurable functions of the sample. It is finite and continuous in \(\gamma\), because
\[
|\widehat I_C(\gamma)-\widehat I_C(\widetilde\gamma)|
\le \sup_{f\in\mathcal F_{\Sigma,s}}\sum_{i\in D_C}|\gamma_i-\widetilde\gamma_i|\,|f(X_i)|.
\]
The feasible box is compact. The positive quadratic term in \(\widehat\Psi_C\) makes the objective strictly convex, hence its minimizer is unique; equivalently it is the lexicographically least minimizer. Choosing successively the least rational-grid point whose objective is within \(n^{-1}\) of the grid infimum gives measurable approximants converging to that unique minimizer. Thus the weights and estimator are jointly measurable and total. The empty-donor, empty-target, stayer, singular-history, and fallback branches were explicitly defined, so no exceptional sample remains.

We next identify the object estimated. Since \(S\) has nonempty full-rank interior, \(Q\) is not the zero polynomial. Its zero set has Lebesgue measure zero; \(\text{ass:mover-density}\) therefore gives \(P\{Q(X)=0\mid X\ne0\}=0\). On \(\{Q(X)>0\}\), \(M(X)^+M(X)=I_p\). By \(\text{ass:distributed-lag-model}\), \(\text{ass:conditional-mean-zero}\), and the \(q\)-moment integrability,
\[
E_P[e_k'M(X)^+(\Delta Y-\Gamma)\mid X]
=E_P[B_k\mid X]=\mu_k(X).
\]
Iterated expectation and the definition of \(\tau_k\) then give
\[
E_P[\mu_k(X)\mid X\ne0]=\tau_k(P).
\]
Thus the causal target is related to, rather than defined as, the observed-law Moore--Penrose functional.

Let \(Z_{gk}^{\circ}=e_k'M(X_g)^+(\Delta Y_g-\Gamma)\) for movers. On the event \(n_V>0\), add and subtract the target-fold empirical mean of \(\mu_k\) to write the oracle-trend estimator error as
\[
\widehat\tau_{k,V,D}^{\circ}(C)-\tau_k
=A_V-B_C(\widehat\gamma_C)+N_V+N_D,
\]
where
\[
A_V=n_V^{-1}\sum_{r\in V:X_r\ne0}\mu_k(X_r)-\tau_k,
\]
\[
B_C(\gamma)=\widehat L_{V,C}(\mu_k)-\sum_{i\in D_C}\gamma_i\mu_k(X_i),
\]
\[
N_V=n_V^{-1}\sum_{r\in V:X_r\ne0}\mathbf 1\{H_k(X_r)\le C\}\{Z_{rk}^{\circ}-\mu_k(X_r)\},
\quad
N_D=\sum_{i\in D_C}\widehat\gamma_{C,i}\{Z_{ik}^{\circ}-\mu_k(X_i)\}.
\]
The mover probability is at least \(1-\overline\pi>0\) by \(\text{ass:stayer-mass}\), so the inverse binomial moment satisfies \(E[\mathbf 1\{n_V>0\}/n_V]\le C/G\), while \(P(n_V=0)\le\overline\pi^{|V|}\). The Whitney bound gives \(\|A_V\|_2\le C G^{-1/2}\). Since \(\mu_k/L_s\in\mathcal F_{\Sigma,s}\),
\[
|B_C(\widehat\gamma_C)|^2
\le L_s^2\widehat I_C(\widehat\gamma_C)^2
\le\widehat\Psi_C(\widehat\gamma_C),
\]
and hence \(\|B_C(\widehat\gamma_C)\|_2\le\Omega_{G,\Sigma,k,s}(C)^{1/2}\).

Conditional on the target histories, \(N_V\) has mean zero. The definition of \(V_{P,k}(C)\) and the conditional identity above imply
\[
E_P[(Z_{k}^{\circ}-\mu_k(X))^2\mathbf 1\{H_k(X)\le C\}\mid X\ne0]
\le V_{P,k}(C),
\]
so \(\|N_V\|_2\le C\{\overline V_{\Sigma,k}(C)/G\}^{1/2}\). For the donor term, independence across groups and \((a+b)^2\le2a^2+2b^2\) give
\[
E_P[N_D^2\mid (X_i)_{i\in D}]
\le 2\sum_{i\in D_C}\widehat\gamma_{C,i}^2E_P[(B_{ik}-\mu_k(X_i))^2\mid X_i]
+2\overline\sigma^2\sum_{i\in D_C}\widehat\gamma_{C,i}^2H_k(X_i)^2.
\]
The cap \(|\widehat\gamma_{C,i}|\le\Lambda_\Sigma/G\), \(\text{ass:q-moments}\), and \(q>4\) bound the expectation of the first sum by \(C/G\). The second sum is at most \(C\widehat\Psi_C(\widehat\gamma_C)\). Consequently
\[
\|N_D\|_2\le C\{G^{-1/2}+\Omega_{G,\Sigma,k,s}(C)^{1/2}\}.
\]

It remains to replace \(\Gamma\) by \(\widehat\Gamma\). By the already-established \(\text{lem:regular-gamma}\), the training-fold error has \(L^2\) norm \(O(G^{-1/2})\) and is independent of donor and target folds. The lower variance bound and the identity
\[
e_k'M(X)^+\varepsilon
=\{Z_k^{\circ}-\tau_k\}-\{B_k-\tau_k\}
\]
give, using \(\text{ass:q-moments}\),
\[
E_P[H_k(X)^2\mathbf 1\{H_k(X)\le C\}\mid X\ne0]
\le C\{1+V_{P,k}(C)\}.
\]
Cauchy--Schwarz therefore bounds the generated-trend error in the direct term by \(C[G^{-1/2}+\{\overline V_{\Sigma,k}(C)/G\}^{1/2}]\). For the donor term,
\[
\left\|\sum_{i\in D_C}\widehat\gamma_{C,i}(M(X_i)')^+e_k\right\|_2^2
\le |D_C|\sum_{i\in D_C}\widehat\gamma_{C,i}^2H_k(X_i)^2,
\]
so independence of the trend fold and the objective penalty bound its generated-trend contribution by \(C\Omega_{G,\Sigma,k,s}(C)^{1/2}\). The exponentially rare projector fallback and empty-target events are absorbed by \(G^{-1/2}\). Averaging a fixed number of rotations and applying the triangle inequality proves the displayed uniform \(L^2\) bound.

Finally, the separate deterministic inequalities follow from \(\text{lem:selector-geometry}\), \(\text{ass:frontier-holder}\), and the density envelope:
\[
\overline{\mathfrak B}_{\Sigma,k,s}(C)
\le C_{\Sigma,s}L_s\mathcal D_{\Sigma,k,s}(C)
\le C'_{\Sigma,s}L_s\mathcal J_{\Sigma,k,s}(C).
\]
No step bounds the empirical optimal-recovery quantity \(\Omega_{G,\Sigma,k,s}(C_{G,\Sigma,k}^{\star})\) by the deterministic pole criterion, so the proof makes no such assertion.
\end{proof}
\par\smallskip\noindent \textit{Justification.} The empirical minimax-linear objective produces a total estimator and exposes its exact optimal-recovery cost. \textit{Gap.} Published local-polynomial and series rates require support, occupancy, or Gram conditions absent from the frozen normal-link class. \textit{Consumer.} The result supplies the attainable augmented-oracle risk certificate for \(\texttt{dist\_lag\_het}\).

\begin{theorem}[thm:matched-frontier]\label{thm:matched-frontier}
For each fixed certified \(\Sigma\), the cross-fitted minimax-linear estimator attains the augmented class oracle: for all sufficiently large \(G\), \[\operatorname{Risk}_{G,k}(\Sigma)\le C R_{G,\Sigma,k}^{\dagger}(C_{G,\Sigma,k}^{\dagger}),\] where \(C<\infty\) depends only on \(\Sigma,\mathbf c,s,q\). Separately, for \(0<\alpha_{\Sigma,k}<2\) write \(v_{\Sigma,k}=2-\alpha_{\Sigma,k}\). The deterministic pole-sensitive oracle obeys \[R_{G,\Sigma,k}^{\star}\lesssim\begin{cases}G^{-1},&2<\alpha_{\Sigma,k}\le\infty,\\G^{-1}(\log G)^{m_{\Sigma,k}+1},&\alpha_{\Sigma,k}=2,\\G^{-2\zeta_{\Sigma,k,s}/(v_{\Sigma,k}+2\zeta_{\Sigma,k,s})}(\log G)^{2\{\zeta_{\Sigma,k,s}m_{\Sigma,k}+\ell_{\Sigma,k,s}v_{\Sigma,k}\}/(v_{\Sigma,k}+2\zeta_{\Sigma,k,s})},&0<\alpha_{\Sigma,k}<2.\end{cases}\] The implicit constants depend only on \(\Sigma,\mathbf c,s,q\). The equality case \(\alpha_{\Sigma,k}=2\), every logarithmic multiplicity, and the bounded-score branch are retained. No comparison \(R_{G,\Sigma,k}^{\dagger}(C_{G,\Sigma,k}^{\dagger})\lesssim R_{G,\Sigma,k}^{\star}\), reverse minimax inequality, or honest-length bound is asserted.
\end{theorem}
\begin{proof}
First apply the proposed oracle-fitting lemma at \(C=C_{G,\Sigma,k}^{\dagger}\). Squaring its \(L^2\) bound and using \((a+b+c)^2\le3(a^2+b^2+c^2)\) gives
\[
\sup_{P\in\mathcal P_\Sigma}E_P[(\widehat\tau_k(C_{G,\Sigma,k}^{\dagger})-\tau_k(P))^2]
\le C\left\{G^{-1}\overline V_{\Sigma,k}(C_{G,\Sigma,k}^{\dagger})
+\Omega_{G,\Sigma,k,s}(C_{G,\Sigma,k}^{\dagger})+G^{-1}\right\}.
\]
The omitted squared-bias-envelope term is nonnegative. By \(\text{def:class-oracle}\), the right-hand side is therefore at most \(C R_{G,\Sigma,k}^{\dagger}(C_{G,\Sigma,k}^{\dagger})\). Since this is one measurable data-driven estimator, the infimum in \(\text{def:minimax-functionals}\) is no larger. This proves the feasible risk claim without an oracle nuisance or a law-dependent estimator.

It remains to derive the separate deterministic phase. On \(\{Q>0\}\), the model and conditional mean restriction give
\[
e_k'M(X)^+(\Delta Y-\Gamma)=B_k+e_k'M(X)^+\varepsilon.
\]
The \(q\)-moment bound, conditional variance upper bound, and \(H_k=\|(M')^+e_k\|_2\) imply
\[
\overline V_{\Sigma,k}(C)
\le C_0\left(1+\int_0^C2u\sup_{P\in\mathcal P_\Sigma}P\{H_k(X)>u\mid X\ne0\}\,du\right).
\]
The tail atlas yields
\[
\overline V_{\Sigma,k}(C)\lesssim
\begin{cases}
1,&\alpha_{\Sigma,k}>2,\\
(\log C)^{m_{\Sigma,k}+1},&\alpha_{\Sigma,k}=2,\\
C^{v_{\Sigma,k}}(\log C)^{m_{\Sigma,k}},&0<\alpha_{\Sigma,k}<2.
\end{cases}
\]
The selector lemma and the density upper bound give
\[
\overline{\mathfrak B}_{\Sigma,k,s}(C)
\lesssim C^{-\zeta_{\Sigma,k,s}}(\log C)^{\ell_{\Sigma,k,s}}.
\]

Because the support has nonempty full-rank interior, there is a fixed \(x_0\in S\) with \(Q(x_0)>0\) and \(H_k(x_0)<\infty\). Hence \(A_C\ne\varnothing\) for all \(C\ge H_k(x_0)\). The corrected grid cap diverges, so every polynomial-logarithmic candidate below is oracle-eligible for all sufficiently large \(G\).

If \(\alpha_{\Sigma,k}>2\), the tail integral gives \(\overline V_{\Sigma,k}(\infty)<\infty\); choosing \(C=\infty\) proves \(R_{G,\Sigma,k}^{\star}\lesssim G^{-1}\). This includes the bounded-score branch. If \(\alpha_{\Sigma,k}=2\), choose a grid cutoff within a fixed factor of a polynomial-logarithmic \(C_G\) satisfying
\[
C_G^{-2\zeta_{\Sigma,k,s}}(\log C_G)^{2\ell_{\Sigma,k,s}}
\le G^{-1}(\log G)^{m_{\Sigma,k}+1}.
\]
The corrected cap in \(\text{def:class-oracle}\) contains such a cutoff for all large \(G\), and the variance bound has the same stated order.

For \(0<\alpha_{\Sigma,k}<2\), abbreviate \(v=v_{\Sigma,k}\), \(\zeta=\zeta_{\Sigma,k,s}\), \(m=m_{\Sigma,k}\), and \(\ell=\ell_{\Sigma,k,s}\). Balancing the two upper bounds gives
\[
C^{v+2\zeta}=G(\log C)^{2\ell-m}.
\]
The grid contains a point within a factor \(a_0\) of
\[
C_G=G^{1/(v+2\zeta)}(\log G)^{(2\ell-m)/(v+2\zeta)}.
\]
Substitution yields
\[
G^{-1}C_G^v(\log C_G)^m
\asymp C_G^{-2\zeta}(\log C_G)^{2\ell}
\asymp
G^{-2\zeta/(v+2\zeta)}
(\log G)^{2(\zeta m+\ell v)/(v+2\zeta)}.
\]
Since \(C_{G,\Sigma,k}^{\star}\) minimizes the deterministic criterion, these choices prove the displayed one-sided phase bound. Neither the fitting lemma nor the normal-link certificate compares \(\Omega_{G,\Sigma,k,s}\) with this deterministic balance, and the frozen class admits the established zero-coefficient obstruction to universal reverse and honest-length bounds. Accordingly the proposed theorem asserts none of those conclusions.
\end{proof}
\par\smallskip\noindent \textit{Justification.} The result separates the proved deterministic Hankel phase from feasible attainment at the augmented empirical-modulus scale. \textit{Gap.} No cited scalar trimming result provides the multivariate resolution-invariant pole phase or closes the empirical frontier-transport modulus. \textit{Consumer.} Users receive distinct geometry and attainable-risk diagnostics without an unsupported honest-length claim.

\begin{proposition}[prop:benchmark-reductions]\label{prop:benchmark-reductions}
The geometry has three reductions. First, let \(P^{\mathrm{fs}}\) be a separate finite-support law under which the groups are independent and identically distributed, the distributed-lag equation and conditional-mean-zero restriction hold, \(0<P^{\mathrm{fs}}(X=0)<1\), and \(E_{P^{\mathrm{fs}}}[\|B\|_2^2+\|\varepsilon\|_2^2]\) is bounded by a fixed finite constant. Suppose that the entire finite mover support is separated from \(\mathcal Z\), equivalently \(Q(x)>0\) for every mover-support point. This discrete law is not an element of the absolutely continuous class \(\mathcal P_\Sigma\). The feasible cross-fitted untrimmed Moore--Penrose estimator using the total positive-stayer fold estimator of \(\Gamma\) identifies \(\tau_k(P^{\mathrm{fs}})\), has squared risk \(O(G^{-1})\), and yields a variance-bound interval with expected length \(O(G^{-1/2})\); the generated-\(\Gamma\) contribution has root mean square \(O(G^{-1/2})\). Second, in the continuous scalar case \(K=0\), \(T=2\), with mover density bounded above and away from zero near zero and \(H_0(x)\asymp|x|^{-1}\), \[\mathcal J_{\Sigma,0,s}(C)\asymp\mathcal D_{\Sigma,0,s}(C)\asymp C^{-(s+1)},\qquad (\alpha_{\Sigma,0},m_{\Sigma,0},\zeta_{\Sigma,0,1},\ell_{\Sigma,0,1})=(1,0,2,0).\] For first-order Whitney extrapolation, \(C\asymp G^{1/5}\) gives the oracle upper bound \(R_{G,\Sigma,0}^{\star}\lesssim G^{-4/5}\); without correction, \(C\asymp G^{1/3}\) gives the upper bound \(G^{-2/3}\). These become two-sided orders when the corresponding uniform bias envelopes are bounded below by positive multiples of \(C^{-2}\) and \(C^{-1}\), respectively. Third, for \(K=1\), \(T=5\) near \(x^\star\), both targets have first leverage pole \((\alpha_{\Sigma,k},m_{\Sigma,k})=(2,0)\), and hence \(R_{G,\Sigma,k}^{\star}\lesssim G^{-1}\log G\); equality requires a nonvanishing frontier modulus that forces the cutoff to diverge.
\end{proposition}
\begin{proof}
Let \(\mathcal X_{\mathrm m}\) denote the finite mover support under \(P^{\mathrm{fs}}\), and put \(\pi_{\mathrm m}=P^{\mathrm{fs}}(X\ne0)\).  The stipulated stayer mass gives \(0<\pi_{\mathrm m}<1\).  The repaired separation premise is
\[
Q(x)>0\qquad\text{for every }x\in\mathcal X_{\mathrm m}.
\]
Because \(Q(x)=\det(M(x)'M(x))\), this is equivalent to full column rank of \(M(x)\) at every mover-support point.  Hence the defining identities of the Moore--Penrose inverse give
\[
M(x)^+M(x)=I_p,
\qquad
e_k'M(x)^+M(x)=e_k',
\qquad x\in\mathcal X_{\mathrm m}.
\]
Finiteness of \(\mathcal X_{\mathrm m}\) then gives the finite leverage bound
\[
\overline H_k^{\mathrm{fs}}
:=\max_{x\in\mathcal X_{\mathrm m}}\lVert e_k'M(x)^+\rVert_2
=\max_{x\in\mathcal X_{\mathrm m}}H_k(x)<\infty.
\]
Define the oracle mover score
\[
S_k=e_k'M(X)^+(\Delta Y-\Gamma).
\]
The distributed-lag equation in the first part of the proposition and its conditional-mean-zero restriction imply, at each \(x\in\mathcal X_{\mathrm m}\),
\[
\begin{aligned}
E_{P^{\mathrm{fs}}}[S_k\mid X=x]
&=E_{P^{\mathrm{fs}}}[e_k'M(x)^+\{M(x)B+\varepsilon\}\mid X=x]\\
&=E_{P^{\mathrm{fs}}}[B_k\mid X=x]
+e_k'M(x)^+E_{P^{\mathrm{fs}}}[\varepsilon\mid X=x]\\
&=E_{P^{\mathrm{fs}}}[B_k\mid X=x].
\end{aligned}
\]
All terms are integrable by the assumed second-moment bound.  Averaging the last equality conditional on \(X\ne0\) yields
\[
E_{P^{\mathrm{fs}}}[S_k\mid X\ne0]
=E_{P^{\mathrm{fs}}}[B_k\mid X\ne0]
=\tau_k(P^{\mathrm{fs}}).
\]
Thus the observed-law score identifies the stated all-mover causal target; equality is derived from full rank and the model, rather than imposed as a definition.

For feasibility, consider a training fold of size \(n\), let \(N_0\) be its number of stayers, and use the following total positive-stayer estimator: if \(N_0>0\), average \(\Delta Y_g\) over the stayers; if \(N_0=0\), use the fold mean of \(\Delta Y_g\).  This fallback is translation equivariant.  At \(X=0\), the definition of \(M\) gives \(M(0)=0\), so the model and conditional-mean-zero condition give
\[
E_{P^{\mathrm{fs}}}[\Delta Y\mid X=0]=\Gamma.
\]
Conditional on \(N_0=m>0\), the stayer average has mean \(\Gamma\) and mean squared error at most \(C/m\).  For \(N_0\sim\operatorname{Binomial}(n,1-\pi_{\mathrm m})\),
\[
E\!\left[\frac{\mathbf 1\{N_0>0\}}{N_0}\right]
\le 2E\!\left[\frac1{N_0+1}\right]
=\frac{2\{1-\pi_{\mathrm m}^{,n+1}\}}{(n+1)(1-\pi_{\mathrm m})}
=O(n^{-1}).
\]
On \(\{N_0=0\}\), all training observations are movers.  Conditional on that event they remain independent draws from the mover law, and \(M(X)B+\varepsilon\) has a finite second moment because the mover support is finite and the stated second moments are finite.  Consequently the fallback contribution is \(O(\pi_{\mathrm m}^{,n})\).  Combining the two cases proves
\[
E_{P^{\mathrm{fs}}}\lVert\widehat\Gamma-\Gamma\rVert_2^2=O(n^{-1}).
\]

On an independent evaluation fold, let \(N_{\mathrm m}\) be the number of movers, average the feasible scores over those movers when \(N_{\mathrm m}>0\), and return zero otherwise.  The oracle score has a finite conditional second moment because
\[
|S_k|^2
\le 2|B_k|^2+2(\overline H_k^{\mathrm{fs}})^2\lVert\varepsilon\rVert_2^2.
\]
The same inverse-binomial calculation, now with success probability \(\pi_{\mathrm m}\), shows that the oracle mover average has mean squared error \(O(n^{-1})\); its empty-mover contribution is \(\tau_k(P^{\mathrm{fs}})^2(1-\pi_{\mathrm m})^n=O(e^{-cn})\).  The score perturbation due to trend estimation is
\[
T_\Gamma=-\left\{N_{\mathrm m}^{-1}
\sum_{g:X_g\ne0}e_k'M(X_g)^+\right\}(\widehat\Gamma-\Gamma)
\]
on \(\{N_{\mathrm m}>0\}\), and is zero otherwise.  Its bracketed row vector has norm at most \(\overline H_k^{\mathrm{fs}}\).  Training/evaluation independence therefore gives
\[
E_{P^{\mathrm{fs}}}[T_\Gamma^2]
\le (\overline H_k^{\mathrm{fs}})^2
E_{P^{\mathrm{fs}}}\lVert\widehat\Gamma-\Gamma\rVert_2^2
=O(G^{-1}).
\]
Using a fixed number of cross-fitting rotations and the inequality \((a+b)^2\le2a^2+2b^2\) proves squared risk \(O(G^{-1})\) and root mean square \(O(G^{-1/2})\) for the generated-\(\Gamma\) term.  If \(C_{\mathrm{fs}}/G\) is a deterministic bound on this mean squared error for all sufficiently large \(G\), then
\[
\mathcal C_{G,k}^{\mathrm{fs}}
=\left[
\widehat\tau_k^{\mathrm{fs}}-\sqrt{C_{\mathrm{fs}}/(\eta G)},
\widehat\tau_k^{\mathrm{fs}}+\sqrt{C_{\mathrm{fs}}/(\eta G)}
\right]
\]
has coverage at least \(1-\eta\) by Markov's inequality applied to squared error and has deterministic, hence expected, length \(2\sqrt{C_{\mathrm{fs}}/(\eta G)}=O(G^{-1/2})\).  This establishes the finite-support reduction under the separate law \(P^{\mathrm{fs}}\), without placing that law in \(\mathcal P_\Sigma\).

Now specialize to \(K=0\) and \(T=2\).  Then \(p=r=d=1\), \(M(x)=(x)\), \(\mathcal Z=\{0\}\), and the global convention for \(H_0\) gives
\[
H_0(x)=|x|^{-1}\quad(x\ne0),
\qquad H_0(0)=0.
\]
For all sufficiently large \(C\), the frontier set is therefore \(0<|x|<C^{-1}\).  The two-sided mover-density bound near zero gives
\[
P\{H_0(X)>C\mid X\ne0\}
=\int_{0<|x|<C^{-1}}f_X(x)\,dx
\asymp C^{-1}.
\]
Because \(\operatorname{dist}(x,\mathcal Z)=|x|\), the definition of the frontier remainder gives
\[
\mathcal J_{\Sigma,0,s}(C)
=\int_{S\cap\{0<|x|<C^{-1}\}}|x|^s\,dx
\asymp C^{-(s+1)}.
\]
For the selector-distance integral, every retained anchor has absolute value at least \(C^{-1}\).  Hence, on \(0<|x|<(2C)^{-1}\), every selector distance is at least \((2C)^{-1}\), and the partition weights sum to one.  This gives the lower bound \(\mathcal D_{\Sigma,0,s}(C)\gtrsim C^{-(s+1)}\).  Conversely, the boundary-compatible scalar chart meeting either incident side contains a retained point at distance at most a signature-dependent multiple of \(C^{-1}\); the prescribed fallback has the same bound.  Thus every selector distance on the frontier is \(O(C^{-1})\), and its frontier has Lebesgue measure \(O(C^{-1})\).  Therefore
\[
\mathcal D_{\Sigma,0,s}(C)\asymp C^{-(s+1)}.
\]
There are no logarithmic ties in this one-dimensional integral.  The definitions of the invariant indices consequently give
\[
(\alpha_{\Sigma,0},m_{\Sigma,0},
\zeta_{\Sigma,0,1},\ell_{\Sigma,0,1})=(1,0,2,0).
\]

On \(\{H_0\le C,X\ne0\}\), the oracle score is \(B_0+\varepsilon/x\).  The upper density, noise-variance, and moment envelopes imply
\[
\overline V_{\Sigma,0}(C)
\lesssim 1+\int_{C^{-1}\le|x|\le\delta_0}x^{-2}\,dx
\lesssim C.
\]
The boundary-extension and selector bounds give
\(\overline{\mathfrak B}_{\Sigma,0,1}(C)\lesssim C^{-2}\).  Choosing a geometric-grid point within a fixed factor of \(G^{1/5}\) therefore yields
\[
R_{G,\Sigma,0}^{\star}
\le G^{-1}\overline V_{\Sigma,0}(C)
+\overline{\mathfrak B}_{\Sigma,0,1}(C)^2
\lesssim G^{-1}C+C^{-4}
\lesssim G^{-4/5}.
\]
Without boundary correction, the omitted-tail bias is at most a constant times its probability, hence is \(O(C^{-1})\).  A cutoff within a fixed factor of \(G^{1/3}\) then gives
\[
G^{-1}\overline V_{\Sigma,0}(C)+O(C^{-2})
\lesssim G^{-2/3}.
\]
For completeness, the variance order is also sharp for the worst-case scale.  Take an admissible scalar law with \(B_0=0\) and centered bounded noise independent of \(X\), with variance in \([\underline\sigma^2,\overline\sigma^2]\).  Such bounded noise respects the moment envelope whenever the scalar class is nonempty: Lyapunov's inequality already forces \(C_q\ge\underline\sigma^q\).  This law belongs to the scalar fixed-signature class and gives
\[
\overline V_{\Sigma,0}(C)
\ge c\int_{C^{-1}\le|x|\le\delta_0}x^{-2}\,dx
\gtrsim C.
\]
If the corrected uniform bias envelope is, as stipulated in the proposition's qualification, bounded below by \(cC^{-2}\) at every effective large cutoff, minimization of \(C/G+C^{-4}\) gives the reverse bound \(R_{G,\Sigma,0}^{\star}\gtrsim G^{-4/5}\).  If the corresponding uncorrected bias envelope is bounded below by \(cC^{-1}\), minimization of \(C/G+C^{-2}\) gives the reverse order \(G^{-2/3}\).  This proves exactly the two conditional two-sided assertions.

Finally, Proposition \(\text{prop:k1t5-early-kill}\) gives \((\alpha_{\Sigma,k},m_{\Sigma,k})=(2,0)\) for both coordinates when \(K=1,T=5\) near \(x^\star\).  Substitution into the critical branch of Theorem \(\text{thm:matched-frontier}\) yields
\[
R_{G,\Sigma,k}^{\star}\lesssim G^{-1}(\log G)^{m_{\Sigma,k}+1}
=G^{-1}\log G.
\]
The tail integration identity
\[
E[H_k(X)^2\mathbf 1\{1<H_k(X)\le C\}\mid X\ne0]
=2\int_1^C uP\{H_k(X)>u\mid X\ne0\}\,du
-C^2P\{H_k(X)>C\mid X\ne0\}
+P\{H_k(X)>1\mid X\ne0\}
\]
and the two-sided \(u^{-2}\) tail show that the critical truncated second moment is of order \(\log C\).  If the frontier bias modulus vanished at some finite cutoff \(C_0\), the oracle criterion at \(C_0\) would be \(O(G^{-1})\), excluding equality of order \(G^{-1}\log G\).  More generally, if that modulus is nonzero at every finite cutoff and tends to zero only as \(C\to\infty\), any oracle sequence whose criterion tends to zero must have \(C_{G,\Sigma,k}^{\star}\to\infty\).  Thus the displayed critical upper bound is unconditional, whereas equality requires precisely the stated nonvanishing frontier modulus forcing a diverging cutoff.
\end{proof}
\par\smallskip\noindent \textit{Justification.} These exact reductions test the distance valuation and ensure that the corrected scalar rate is not confused with the uncorrected slow-mover balance. \textit{Gap.} The finite-support and uncorrected scalar endpoints appear separately in deChaisemartinDHa2025 and GrahamPowell2012, but neither is embedded in a target-specific Hankel phase diagram with distance-weighted correction. \textit{Consumer.} The reductions give \(\texttt{dist\_lag\_het}\) backward-compatible finite-support behavior when every mover-support point has \(Q(x)>0\), together with a checkable scalar unit benchmark.

\begin{lemma}[lem:zero-coefficient-certified-witness]\label{lem:zero-coefficient-certified-witness}
There exists a nonempty certified fixed-signature class with \(T=2\), \(K=0\), and \(S=[2,3]\) for which \(\tau_0(P)=0\) for every member, \(\operatorname{Risk}_{G,0}(\Sigma)=\operatorname{Len}_{G,0}(\Sigma,\eta)=0\), but \(R_{G,\Sigma,0}^{\star}=\sigma^2/(6G)>0\). Thus nonemptiness and the frozen envelope restrictions do not provide the interior coefficient-mean richness needed for a minimax or honest-length lower bound in terms of \(R_{G,\Sigma,k}^{\star}\).
\end{lemma}
\begin{proof}
Fix \(q>4\), a positive integer \(G\), and \(\eta\in(0,1/2)\). Choose \(\pi\in(0,1)\) and \(\sigma>0\), and specialize the envelope constants to
\[
 \underline\pi=\overline\pi=\pi,
 \qquad
 \underline f=\overline f=1,
 \qquad
 \underline\sigma=\overline\sigma=\sigma,
 \qquad
 C_q=(1-\pi)\sigma^q,
\]
with any fixed \(L_s>0\). Take \(T=2\), \(K=0\), hence \(p=r=d=1\), and take
\[
 S=[2,3],\qquad \mathbf g=(x-2,3-x).
\]
By the definition of \(M\), \(M(x)=(x)\). Consequently
\[
 \mathcal I_{\rm rk}=\langle x\rangle,
 \qquad Q(x)=x^2,
 \qquad A_0(x)=1,
 \qquad \mathcal Z=\{0\}.
\]
This support admits the certificate required by \(\text{ass:normal-link-certificate}\): use the identity chart on the interval, stratify it into \((2,3)\), \(\{2\}\), and \(\{3\}\), and use the inward affine rays at the two endpoint strata. On radius \(\delta_0=1/4\), every nonempty zero-dimensional spherical link has Hausdorff mass one, while all affine Lipschitz, angle, and transverse-Jacobian constants equal one. Thus, for example, the strict rational bounds \(1/2<1<3/2\) certify every nonempty link, and the polynomial sign conditions \(x-2\ge0\) and \(3-x\ge0\) certify the active faces. Rank loss and alignment loss do not meet \(S\), since \(Q\ge4\) and \(A_0=1\) there. These explicit rational data form a finite normal-link record, and hence a certified signature \(\Sigma\) satisfying \(\text{ass:fixed-signature-support}\) and \(\text{ass:normal-link-certificate}\).

The resulting class is nonempty. Indeed, let \(P(X=0)=\pi\), let \(X\mid\{X\ne0\}\) be uniform on \([2,3]\), set \(B=0\) and \(\Gamma=0\), and let
\[
 \varepsilon=0\quad\text{on }\{X=0\},
 \qquad
 P(\varepsilon=\sigma\mid X=x)
 =P(\varepsilon=-\sigma\mid X=x)=\frac12
 \quad\text{for a.e. }x\in[2,3].
\]
Put \(\Delta Y=\varepsilon\), and sample \((W_g)_{g=1}^G\) independently from this law. Then \(\Delta Y=\Gamma+M(X)B+\varepsilon\), so \(\text{ass:distributed-lag-model}\) holds, and the symmetric conditional law gives \(E_P[\varepsilon\mid X]=0\), as required by \(\text{ass:conditional-mean-zero}\). Moreover,
\[
 E_P[|B|^q+|\varepsilon|^q]=(1-\pi)\sigma^q=C_q,
 \qquad
 E_P[\varepsilon^2\mid X=x]=\sigma^2
 \quad\text{for a.e. }x\in S.
\]
Thus \(\text{ass:q-moments}\) and \(\text{ass:noise-variance}\) hold. The mover density and stayer mass equal their specified bounds, and \(\mu_0=0\) has the zero \(C^{s-1,1}\) Whitney extension, so all remaining membership conditions in \(\text{def:fixed-signature-class}\) also hold.

We next show that every law in this same fixed-envelope class has zero coefficient. Let \(P\) be any member. The equal stayer bounds imply \(P(X\ne0)=1-\pi\). For mover-almost every \(x\), put \(m_x=E_P[\varepsilon^2\mid X=x]\). By \(\text{ass:noise-variance}\), \(m_x\ge\sigma^2>0\). We derive, rather than assume, the needed conditional moment inequality. For \(a=q/2>1\), the elementary inequality
\[
 t^a\ge 1+a(t-1),\qquad t\ge0,
\]
follows because the derivative of \(t^a-1-a(t-1)\) is \(a(t^{a-1}-1)\), which is nonpositive on \([0,1]\), nonnegative on \([1,\infty)\), and the function vanishes at one. Apply this inequality to \(t=\varepsilon^2/m_x\) and take the conditional expectation to obtain
\[
 E_P[|\varepsilon|^q\mid X=x]
 =m_x^{q/2}E_P[(\varepsilon^2/m_x)^{q/2}\mid X=x]
 \ge m_x^{q/2}
 \ge\sigma^q.
\]
Integration over the mover event therefore gives
\[
 E_P|\varepsilon|^q\ge(1-\pi)\sigma^q.
\]
Using \(\text{ass:q-moments}\) and the chosen value of \(C_q\),
\[
 (1-\pi)\sigma^q=C_q
 \ge E_P[|B|^q+|\varepsilon|^q]
 \ge E_P|B|^q+(1-\pi)\sigma^q.
\]
Hence \(E_P|B|^q=0\), so \(B=0\) almost surely. Since \(1-\pi>0\), the causal target is well defined and
\[
 \tau_0(P)=E_P[B\mid X\ne0]=0
\]
for every member \(P\) of the class.

It follows directly from \(\text{def:minimax-functionals}\) that the constant estimator zero has worst-case squared error zero. Nonnegativity of squared error gives
\[
 \operatorname{Risk}_{G,0}(\Sigma)=0.
\]
Likewise, the deterministic interval \([0,0]\) covers \(\tau_0(P)\) with probability one for every member and has length zero. Since interval length is nonnegative,
\[
 \operatorname{Len}_{G,0}(\Sigma,\eta)=0.
\]

It remains to compute the deterministic oracle, which is an observed-score functional and is distinct from the causal target. For \(x\in S\), the scalar Moore--Penrose inverse is \(M(x)^+=x^{-1}\); therefore, by the definition of \(H_0\),
\[
 H_0(x)=x^{-1}\le\frac12.
\]
Every finite cutoff in \(\mathcal G_{G,\Sigma}\) is at least one, and the grid also contains infinity. Thus \(\mathbf 1\{H_0(X)\le C\}=1\) on the mover event for every oracle-eligible \(C\). The distributed-lag equation and \(B=0\) give
\[
 e_0'M(X)^+(\Delta Y-\Gamma)-\tau_0
 =\frac{\varepsilon}{X}.
\]
Because the equal lower and upper variance bounds force \(E_P[\varepsilon^2\mid X=x]=\sigma^2\), while the equal density bounds force \(f_X(x)=1\) almost everywhere on \([2,3]\), the variance definition in \(\text{def:uniform-scales}\) yields, for every member \(P\) and every oracle-eligible \(C\),
\[
 V_{P,0}(C)
 =\int_2^3\frac{\sigma^2}{x^2}\,dx
 =\sigma^2\left(\frac12-\frac13\right)
 =\frac{\sigma^2}{6}.
\]
Consequently \(\overline V_{\Sigma,0}(C)=\sigma^2/6\). Also \(\{x\in S:H_0(x)>C\}=\varnothing\) for every finite eligible \(C\), so the definition of the uniform bias envelope gives \(\overline{\mathfrak B}_{\Sigma,0,s}(C)=0\); it is zero at infinity by definition as well. Hence every eligible cutoff has the same objective in \(\text{def:class-oracle}\), and
\[
 R_{G,\Sigma,0}^{\star}
 =G^{-1}\frac{\sigma^2}{6}+0
 =\frac{\sigma^2}{6G}>0.
\]
This proves all assertions of the witness.
\end{proof}

\begin{theorem}[thm:no-universal-fixed-signature-lower-bound]\label{thm:no-universal-fixed-signature-lower-bound}
The answer is negative for the frozen class quantified over every nonempty certified signature. There is no universal \(c>0\) for which both \(\operatorname{Risk}_{G,k}(\Sigma)\ge cR_{G,\Sigma,k}^{\star}\) and \(\operatorname{Len}_{G,k}(\Sigma,\eta)\ge c\sqrt{R_{G,\Sigma,k}^{\star}}\) hold for all such classes. A necessary repair is a quantitative interior-slack and coefficient-mean-richness condition permitting two-sided target perturbations while preserving the \(q\)-moment and variance envelopes.
\end{theorem}
\begin{proof}
Let \(c>0\) be arbitrary. Apply \(\text{lem:zero-coefficient-certified-witness}\) to its nonempty certified signature with \(T=2\), \(K=0\), and \(k=0\). That lemma gives
\[
 \operatorname{Risk}_{G,0}(\Sigma)=0,
 \qquad
 \operatorname{Len}_{G,0}(\Sigma,\eta)=0,
 \qquad
 R_{G,\Sigma,0}^{\star}=\frac{\sigma^2}{6G}>0.
\]
Therefore the proposed universal risk inequality would imply
\[
 0=\operatorname{Risk}_{G,0}(\Sigma)
 \ge cR_{G,\Sigma,0}^{\star}
 =\frac{c\sigma^2}{6G}>0,
\]
a contradiction. The proposed length inequality would separately imply
\[
 0=\operatorname{Len}_{G,0}(\Sigma,\eta)
 \ge c\sqrt{R_{G,\Sigma,0}^{\star}}
 =\frac{c\sigma}{\sqrt{6G}}>0,
\]
which is also a contradiction. Since \(c>0\) was arbitrary, no positive universal constant can make either inequality, and hence cannot make both inequalities, hold over every nonempty certified fixed-signature class.

The same witness identifies the obstruction rather than merely refuting the inequalities. Its variance lower bound requires mover noise with conditional \(q\)-moment at least \(\sigma^q\), and its unconditional \(q\)-moment budget is exactly \((1-\pi)\sigma^q\). The equation-by-equation calculation in \(\text{lem:zero-coefficient-certified-witness}\) therefore forces \(E_P|B|^q=0\) for every member. In particular, the class contains no two laws whose coefficient means, and hence whose targets, differ. Any two-sided local perturbation of the coefficient mean would add a strictly positive amount to \(E_P|B|^q\) while the mandatory noise already exhausts the budget, so it would leave the class. Thus a lower-bound repair that excludes this obstruction must impose quantitative slack relative to the moment and variance envelopes and quantitative richness permitting target-changing coefficient-mean perturbations in both directions. This proves the stated negative answer and its necessary repair.
\end{proof}

\begin{theorem}[thm:no-universal-joint-adaptation-factor]\label{thm:no-universal-joint-adaptation-factor}
The requested universal matching adaptation factor does not exist under the frozen class. Even with a fixed certified \(\Sigma\), there is no positive constant factor lower bound in units of \(R_{G,\Sigma,k}^{\star}\) or \(\sqrt{R_{G,\Sigma,k}^{\star}}\) that holds over every nonempty class. Moreover, the four resolution indices are deterministic functions of the allowed input \(\Sigma\), so adaptation to them is not a separate statistical experiment within one fixed signature. A positive joint theorem requires a quantitatively rich union of signatures or submodels and a fully specified comparison statistic and critical value.
\end{theorem}
\begin{proof}
First consider any asserted positive factor lower bound, uniform over all nonempty certified fixed-signature classes, in units of \(R_{G,\Sigma,k}^{\star}\) for squared risk or \(\sqrt{R_{G,\Sigma,k}^{\star}}\) for expected confidence length. By \(\text{lem:zero-coefficient-certified-witness}\), there is a nonempty certified class for which
\[
 \operatorname{Risk}_{G,0}(\Sigma)
 =\operatorname{Len}_{G,0}(\Sigma,\eta)=0,
 \qquad
 R_{G,\Sigma,0}^{\star}=\frac{\sigma^2}{6G}>0.
\]
Multiplication of the positive oracle scale, or its square root, by any positive factor remains strictly positive. Hence the displayed equalities contradict either type of uniform positive-factor lower bound. This proves the first assertion.

For the adaptation assertion, \(\text{def:invariant-indices}\) computes
\[
 I_k(\Sigma)
 =\bigl(\alpha_{\Sigma,k},m_{\Sigma,k},
          \zeta_{\Sigma,k,s},\ell_{\Sigma,k,s}\bigr)
\]
from the valuations and incident faces of the certified atlas. The certified atlas is itself fixed by \(\Sigma\). Thus, if \(P_1,P_2\in\mathcal P_\Sigma\), then
\[
 I_k(\Sigma;P_1)=I_k(\Sigma)=I_k(\Sigma;P_2).
\]
In particular, the procedure receives \(\Sigma\) and therefore receives enough deterministic input to compute the same four indices before observing \(W_1,\ldots,W_G\). There are no two laws inside one fixed-signature experiment that differ in these indices, so learning or adapting to them is not a distinct within-signature statistical problem. To formulate a nonvacuous joint adaptation experiment, the parameter space must instead contain a quantitatively rich union of signatures, or explicitly indexed submodels, across which the relevant benchmark data can differ; the zero-coefficient witness shows why nonemptiness alone is not the required richness.

Finally, \(\text{def:adaptive-handle}\) lists the data features to be used and names a Goldenshluger--Lepski selector, but it supplies no formula for a pairwise comparison statistic, no threshold or penalty defining acceptance, and no critical value or radius defining \(\mathcal C_{G,k}\). Therefore that construction does not determine a unique measurable map
\[
 (W_1,\ldots,W_G)\longmapsto
 (\widehat C_k,\mathcal C_{G,k}).
\]
Without such a map, neither the distribution of the selected cutoff nor the coverage probability and expected length of the interval are defined by the frozen construction, so no positive uniform adaptation-and-coverage theorem for that handle can be proved. A positive joint theorem consequently requires both a rich comparison experiment and a fully specified comparison statistic, threshold, and interval critical value, as claimed.
\end{proof}

\begin{lemma}[lem:published-reduced-form-scope]\label{lem:published-reduced-form-scope}
De Chaisemartin and D'Haultf\oe{}uille's Theorem 4 is stated under the following potential-outcome, sampling, moment, and design hypotheses. Write \(\mathbf 0_t\) for the length-\(t\) zero treatment path and \(D_g=(D_{g,1},\ldots,D_{g,T})\) for group \(g\)'s treatment path. First, no anticipation requires that, for every group \(g\), period \(t<T\), and supported treatment path \((d_1,\ldots,d_T)\), \(Y_{g,t}(d_1,\ldots,d_T)\) does not depend on \((d_{t+1},\ldots,d_T)\). Second, the group-indexed collection consisting of the entire potential-outcome schedule and \(D_g\) is independent and identically distributed across groups. Third, parallel trends for the never-treated outcome requires, for every \(t\ge2\),
\[
E\!\left[Y_t(\mathbf 0_t)-Y_{t-1}(\mathbf 0_{t-1})\mid D\right]
=E\!\left[Y_t(\mathbf 0_t)-Y_{t-1}(\mathbf 0_{t-1})\right].
\]
Fourth, for a known \(K<T-1\), there is a random vector \(B=(B_0,\ldots,B_K)'\) with \(E[\lVert B\rVert_2^2]<\infty\) such that
\[
Y_t(d_1,\ldots,d_t)
=Y_t(\mathbf 0_t)+\sum_{j=0}^{\min\{K,t-1\}}B_jd_{t-j}.
\]
Fifth, the design condition requires \(E[\Pi(M)]\) to be invertible. Finally, the additional moment condition is
\[
E\!\left[(D_t-D_{t-1})^2+\{Y_t(\mathbf 0_t)-Y_{t-1}(\mathbf 0_{t-1})\}^2\right]<\infty
\qquad(2\le t\le T).
\]
These potential-outcome and sampling hypotheses imply the induced reduced-form relations
\[
\Delta Y=\Gamma+MB+\varepsilon,
\qquad E[\varepsilon\mid D]=0,
\qquad E[\varepsilon\mid M]=0.
\]
They are logically additional to those reduced-form relations. Under the full displayed hypotheses, Theorem 4 identifies
\[
\Gamma=E[\Pi(M)]^{-1}E[\Pi(M)\Delta Y].
\]
For every coordinate \(k\) with \(P\{e_k\in\operatorname{Im}(M')\}>0\), it defines the component-specific eligible mean
\[
\overline\beta_k=E[B_k\mid e_k\in\operatorname{Im}(M')]
\]
and identifies it by
\[
\overline\beta_k
=\lim_{C\uparrow\infty}
E\!\left[e_k'M^+(\Delta Y-\Gamma)
\mathrel{\big|}e_k\in\operatorname{Im}(M'),\ \lVert(M')^+e_k\rVert_2\le C\right],
\]
where the conditional expectations are taken at cutoffs having positive conditioning probability. Eligibility is coordinate specific. The discussion following Theorem 4 gives the \(K=1\), \(T=5\) support \(\{(0,1,0,1,0),(0,0,1,1,0)\}\), for which the design condition holds and both coordinate means are identified although \(P(D_1=D_2=D_3)=0\). More generally, it explains that overidentification can permit the projector condition without a sufficiently long stayer when \(T>2(K+1)\).
\end{lemma}

\begin{proposition}[prop:published-source-bridge]\label{prop:published-source-bridge}
Fix \(k\in\{0,\ldots,K\}\). For every \(P\in\mathcal P_\Sigma(T,K,s,q,\mathbf c)\),
\[
P\{Q(X)>0\mid X\ne0\}=1,
\qquad
\{e_k\in\operatorname{Im}(M(X)')\}=\{X\ne0\}
\quad P\text{-almost surely},
\]
and
\[
E_P[\Pi(M(X))]\succeq \pi_0 I_r\succ0.
\]
Moreover, \(K<T-1\), \(E_P[\lVert B\rVert_2^2]<\infty\), and \(E_P[\varepsilon\mid M(X)]=0\). Thus the induced reduced-form law satisfies de Chaisemartin--D'Haultf\oe{}uille's distributed-lag design condition, and its component-specific eligible mean obeys
\[
\overline\beta_k
:=E_P[B_k\mid e_k\in\operatorname{Im}(M(X)')]
=E_P[B_k\mid X\ne0]
=\tau_k.
\]
For finite \(C\), let
\[
\rho_{k,C}:=P\{H_k(X)\le C\mid X\ne0\}.
\]
Whenever \(\rho_{k,C}>0\), this note's indicator/mover-denominator functional and the published conditional-cutoff functional satisfy
\[
E_P\!\left[e_k'M(X)^+(\Delta Y-\Gamma)\mathbf 1\{H_k(X)\le C\}\mid X\ne0\right]
=\rho_{k,C}
E_P\!\left[e_k'M(X)^+(\Delta Y-\Gamma)
\mathrel{\big|}e_k\in\operatorname{Im}(M(X)'),\ H_k(X)\le C\right].
\]
The two finite-\(C\) normalizations need not be equal, but both limits as \(C\to\infty\) equal \(\overline\beta_k=\tau_k\). Consequently, \(\mathcal P_\Sigma\) is a strict compact-semialgebraic, full-rank-mover, positive-stayer, bounded-density, bounded-moment, bounded-variance, Whitney-smooth specialization of the published induced reduced-form experiment. It is not the paper's published potential-outcome class. For its lag dimension and reduced-form design requirements, the latter experiment uses only \(K<T-1\) and invertibility of \(E[\Pi(M)]\); it permits component-specific eligibility and overidentified designs without a stayer atom.
\end{proposition}
\begin{proof}
Fix \(P\in\mathcal P_\Sigma(T,K,s,q,\mathbf c)\) and \(k\in\{0,\ldots,K\}\). The upper bound \(\pi_0\le\overline\pi<1\) in \(\text{ass:stayer-mass}\) gives \(P(X\ne0)=1-\pi_0>0\), so every conditional mover expectation below is well defined.

First we prove that rank loss is mover-null. By \(\text{ass:at-least-square}\),
\[
r=T-K-1\ge K+1=p.
\]
Define \(x^*\in\mathbb R^d\) by \(x^*_{K+1}=1\) and \(x^*_a=0\) for \(a\ne K+1\). This coordinate exists because \(d=T-1\ge K+1\). For \(1\le i\le p\) and \(0\le j\le K\), the definition of the Hankel map gives
\[
M(x^*)_{i,j+1}=x^*_{K+i-j}=\mathbf 1\{i=j+1\}.
\]
Thus the first \(p\) rows of \(M(x^*)\) form \(I_p\), so \(M(x^*)\) has column rank \(p\) and
\[
Q(x^*)=\det\{M(x^*)'M(x^*)\}>0.
\]
In particular, \(Q\) is not the zero polynomial.

For completeness, a nonzero real polynomial on \(\mathbb R^d\) has a Lebesgue-null zero set. In dimension one it has finitely many roots. Inductively, write the polynomial as a polynomial in its last coordinate, with coefficient polynomials in the first \(d-1\) coordinates. By the induction hypothesis, the set where all nonzero coefficient polynomials vanish is null. Away from that set, every one-dimensional section has finitely many roots, so Fubini's theorem makes the full zero set null. Applying this fact to \(Q\), and using the conditional mover density in \(\text{ass:mover-density}\), yields
\[
P\{Q(X)=0\mid X\ne0\}
=\int_{S\cap\{Q=0\}}f_X(x)\,dx=0.
\]
Since \(Q\ge0\),
\[
P\{Q(X)>0\mid X\ne0\}=1.
\]

If \(Q(x)>0\), then \(M(x)'M(x)\) is nonsingular and \(M(x)\) has column rank \(p\). Hence \(\operatorname{Im}(M(x)')=\mathbb R^p\), so \(e_k\in\operatorname{Im}(M(x)')\). At \(x=0\), the definition of \(M\) gives \(M(0)=0\), whose row space is \(\{0\}\); because \(e_k\ne0\), it is not eligible. The only remaining histories are nonzero singular movers, which form a null event by the preceding display. Therefore
\[
\{e_k\in\operatorname{Im}(M(X)')\}
=\{X\ne0\}
\quad P\text{-almost surely}.
\]

Next, \(\Pi(M)=I_r-MM^+\) is an orthogonal projector and hence is positive semidefinite. Since \(M(0)=0\), \(\Pi(M(0))=I_r\). The stayer-mass assumption therefore gives
\[
E_P[\Pi(M(X))]
=\pi_0I_r+E_P[\Pi(M(X))\mathbf 1\{X\ne0\}]
\succeq\pi_0I_r
\succeq\underline\pi I_r
\succ0.
\]
Thus \(E_P[\Pi(M(X))]\) is invertible. Also, \(T\ge2K+2\) implies \(K<T-1\). Since \(q>4\), H\"older's inequality and \(\text{ass:q-moments}\) give
\[
E_P[\lVert B\rVert_2^2]
\le \{E_P[\lVert B\rVert_2^q]\}^{2/q}
\le C_q^{2/q}<\infty.
\]
Because \(M(X)\) is measurable with respect to \(X\), the tower property and \(\text{ass:conditional-mean-zero}\) give
\[
E_P[\varepsilon\mid M(X)]
=E_P\!\left[E_P[\varepsilon\mid X]\mid M(X)\right]=0.
\]
Together with the reduced-form equation in \(\text{ass:distributed-lag-model}\) and independent sampling in \(\text{ass:iid-groups}\), these statements verify the reduced-form dimension, moment, exclusion, and projector-invertibility requirements recorded in \(\text{lem:published-reduced-form-scope}\). This is an observed-law bridge only: \(\text{def:fixed-signature-class}\) does not impose the source paper's logically additional no-anticipation and never-treated parallel-trends restrictions on an entire potential-outcome schedule.

The common trend is itself an observed-law functional. Indeed, \(\Pi(M)M=0\) pointwise and the preceding conditional-mean-zero calculation also gives \(E_P[\Pi(M(X))\varepsilon]=0\). Applying \(\Pi(M(X))\) to \(\text{ass:distributed-lag-model}\), taking expectations, and using that \(\Gamma\) is deterministic yields
\[
E_P[\Pi(M(X))\Delta Y]
=E_P[\Pi(M(X))]\Gamma.
\]
The proved strict positivity therefore implies the explicit identifying map
\[
\Gamma
=E_P[\Pi(M(X))]^{-1}E_P[\Pi(M(X))\Delta Y].
\]

The almost-sure event equality now yields
\[
\overline\beta_k
=E_P[B_k\mid e_k\in\operatorname{Im}(M(X)')]
=E_P[B_k\mid X\ne0]
=\tau_k.
\]
This display proves the relationship between the induced reduced-form eligible mean and the causal mover target; it does not create that relationship by definition.

It remains to compare the two cutoff normalizations. Put
\[
R_k=e_k'M(X)^+(\Delta Y-\Gamma),
\qquad
E_C=\{X\ne0,\ H_k(X)\le C\}.
\]
On \(E_C\), full column rank holds almost surely, so \(M(X)^+M(X)=I_p\). The distributed-lag equation gives
\[
R_k=B_k+e_k'M(X)^+\varepsilon
\quad\text{almost surely on }E_C.
\]
For finite \(C\), the error term is integrable on \(E_C\): the identity \((M')^+=(M^+)'\), Cauchy--Schwarz in \(\mathbb R^r\), and the definition of \(H_k\) give
\[
E_P\!\left[|e_k'M(X)^+\varepsilon|\mathbf 1_{E_C}\right]
\le C E_P[\lVert\varepsilon\rVert_2]<\infty,
\]
where the last inequality follows from \(\text{ass:q-moments}\). The Moore--Penrose map is continuous on each fixed-rank matrix stratum; the finitely many rank strata are Borel because rank is characterized by minors. Hence \(M(X)^+\), \(E_C\), and \(e_k'M(X)^+\) are measurable functions of \(X\). Conditional mean zero then implies
\[
E_P[e_k'M(X)^+\varepsilon\mathbf 1_{E_C}]=0.
\]
Consequently,
\[
E_P[R_k\mathbf 1\{H_k(X)\le C\}\mid X\ne0]
=E_P[B_k\mathbf 1\{H_k(X)\le C\}\mid X\ne0].
\]
Whenever
\[
\rho_{k,C}=P(E_C\mid X\ne0)>0,
\]
the eligibility-event equality likewise gives
\[
E_P[R_k\mid e_k\in\operatorname{Im}(M(X)'),\ H_k(X)\le C]
=\frac{E_P[B_k\mathbf 1\{H_k(X)\le C\}\mid X\ne0]}{\rho_{k,C}}.
\]
Multiplying the last display by \(\rho_{k,C}\) proves the asserted finite-cutoff relation.

The Moore--Penrose inverse of every finite real matrix is finite, so \(H_k(X)<\infty\) almost surely. Hence \(\rho_{k,C}\uparrow1\) as \(C\to\infty\). The moment condition makes \(B_k\) integrable, and dominated convergence gives
\[
E_P[B_k\mathbf 1\{H_k(X)\le C\}\mid X\ne0]
\longrightarrow E_P[B_k\mid X\ne0]=\tau_k.
\]
The indicator/mover-denominator functional therefore converges to \(\tau_k\), while division by \(\rho_{k,C}\to1\) shows that the conditional-cutoff functional has the same limit. Thus only the limiting functionals are equal in general.

Finally, membership in \(\mathcal P_\Sigma\) imposes, by \(\text{def:fixed-signature-class}\), a fixed compact full-dimensional semialgebraic support and face signature, mover absolute continuity with a two-sided density envelope, positive stayer mass, \(q\)-moment and conditional-variance envelopes, a normal-link certificate, and original-coordinate Whitney smoothness. The preceding argument adds full-rank movers almost surely. These restrictions make \(\mathcal P_\Sigma\) a specialization of the published induced reduced-form experiment.

The specialization is strict. Retain the same \(T\) and \(K\), take \(X=0\) and \(X=x^*\) with probabilities one half each, set \(D_1=0\) and \(D_{t+1}=D_t+X_t\), and define every potential outcome to be zero. Equivalently, set \(B=0\), \(\varepsilon=0\), \(\Gamma=0\), and \(\Delta Y=0\), with independent copies of the full treatment and potential-outcome schedule across groups. Then
\[
E[\Pi(M(X))]
=\tfrac12I_r+\tfrac12\Pi(M(x^*))
\succeq\tfrac12I_r\succ0.
\]
The reduced-form equation, exclusion, finite-moment, dimension, independent-sampling, and positive-eligibility requirements hold, but the conditional mover distribution is the point mass at \(x^*\). It therefore violates \(\text{ass:mover-density}\) and belongs to no absolutely continuous fixed-signature class of the displayed form. This explicit witness proves strictness. The no-stayer overidentified example and the component-specific eligibility scope recorded in \(\text{lem:published-reduced-form-scope}\) additionally show why neither the fixed atlas nor any theorem depending on it is being extended to the unrestricted published experiment. This completes the proof.
\end{proof}
\par\smallskip\noindent \textit{Justification.} This proposition supplies the exact observed-law bridge: the fixed-signature class is a strict specialization of the anchor's induced reduced-form experiment, and the causal mover mean equals the anchor's eligible mean only after the event equality is proved. \textit{Gap.} The anchor's Theorem 4 uses additional potential-outcome hypotheses and permits component-specific eligibility and overidentified designs without stayers, whereas this note's tail theory requires a positive stayer atom and a fixed absolutely continuous semialgebraic mover geometry. \textit{Consumer.} The bridge gives the \(\texttt{dist\_lag\_het}\) consumer the precise common target while preventing the fixed-signature tail atlas from being stated over either the anchor's unrestricted induced reduced-form experiment or its published potential-outcome class.

\section*{Technical internal limitation (diagnostic only)}

The leverage-tail pair alone does not determine post-correction bias. The distance-weighted local-zeta pair is separate, and the selector-distance comparison is therefore an explicit lemma rather than an assumption. The absolutely continuous class \(\mathcal P_\Sigma\) does not contain the separate finite-support comparison law and is only an induced reduced-form specialization of the anchor; the note does not verify the anchor's full potential-outcome hypotheses for every law in \(\mathcal P_\Sigma\). Exact minimax equivalence is not claimed without a baseline law and same-signature Whitney perturbations jointly realizing both orders. The fixed-signature upper bound also does not imply adaptation across signatures, and coordinatewise shortest intervals do not determine shortest simultaneous-region volume when several lags share an active stratum.

\section{Honest scope}

Closed theorem targets are the source-faithful record of the anchor's full Theorem 4 hypotheses, the observed-law bridge equating the mover target and eligible mean on \(\mathcal P_\Sigma\), the effective fixed-signature certificate, resolution-invariant local-zeta tail and distance-weighted laws on the finite-exponent branch, eventual zero tails on the bounded-score branch, Borel anchor selection, the selector-distance comparison, total minimax-linear fitting, the class-level augmented-oracle risk upper bound, and the benchmark reductions. The bridge does not claim that \(\mathcal P_\Sigma\) is a subclass of the published potential-outcome model. Exact minimax-risk equivalence remains open until a nonempty baseline and same-signature coefficient-mean perturbations are constructed with all Whitney, moment, variance, and stayer restrictions. Honest-length rates or upper bounds over \(\mathcal P_\Sigma\) remain open; the \(O(G^{-1/2})\) interval belongs only to the separate finite-support benchmark \(P^{\mathrm{fs}}\). Fully data-measurable selection of the dominating geometry and simultaneous lag inference is also open. The \(K=1,T=5\) early-kill statement retains the prespecified alternating history and both target coordinates; the twenty-million-draw uniform-box diagnostic at radius \(0.15\) placed \(u^2P(H_k>u)\) between \(20.10\) and \(20.30\) for \(65\le u\le175\), consistent with the required finite positive limit.

\begin{thebibliography}{99}

  \bibitem{deChaisemartinDHa2025} Clément de Chaisemartin and Xavier D'Haultfœuille (2025), Treatment-Effect Estimation in Complex Designs under a Parallel-Trends Assumption, arXiv:2508.07808.

  \bibitem{GrahamPowell2012} Bryan S. Graham and James L. Powell (2012), Identification and Estimation of Average Partial Effects in ``Irregular'' Correlated Random Coefficient Panel Data Models, Econometrica 80(5), 2105--2152, DOI: 10.3982/ECTA8220.

  \bibitem{KhanTamer2010} Shakeeb Khan and Elie Tamer (2010), Irregular Identification, Support Conditions, and Inverse Weight Estimation, Econometrica 78(6), 2021--2042, DOI: 10.3982/ECTA7372.

  \bibitem{ArellanoBonhomme2012} Manuel Arellano and Stéphane Bonhomme (2012), Identifying Distributional Characteristics in Random Coefficients Panel Data Models, Review of Economic Studies 79(3), 987--1020, DOI: 10.1093/restud/rdr045.

  \bibitem{ConcaEtAl2018} Aldo Conca, Maral Mostafazadehfard, Anurag K. Singh, and Matteo Varbaro (2018), Hankel Determinantal Rings Have Rational Singularities, Advances in Mathematics 335, 111--129, DOI: 10.1016/j.aim.2018.06.011.

  \bibitem{Wongkew1993} Richard Wongkew (1993), Volumes of Tubular Neighbourhoods of Real Algebraic Varieties, Pacific Journal of Mathematics 159(1), 177--184, DOI: 10.2140/pjm.1993.159.177.

  \bibitem{PhongSturm2000} D. H. Phong and Jacob Sturm (2000), Algebraic Estimates, Stability of Local Zeta Functions, and Uniform Estimates for Distribution Functions, Annals of Mathematics 152(1), 277--329, DOI: 10.2307/2661384.

  \bibitem{DinhPham2018} Si Tiep Dinh and Tien Son Pham (2018), Łojasiewicz Inequalities with Explicit Exponents for Smallest Singular Value Functions, Journal of Complexity 44, 44--67, DOI: 10.1016/j.jco.2016.11.007.

  \bibitem{SasakiUra2022} Yuya Sasaki and Takuya Ura (2022), Estimation and Inference for Moments of Ratios with Robustness against Large Trimming Bias, Econometric Theory 38(1), 66--112, DOI: 10.1017/S0266466621000025.

  \bibitem{SasakiUra2026} Yuya Sasaki and Takuya Ura (2026), Slow Movers in Panel Data, Econometric Theory, DOI: 10.1017/S0266466625100157.

  \bibitem{Donoho1994} David L. Donoho (1994), Statistical Estimation and Optimal Recovery, Annals of Statistics 22(1), 238--270, DOI: 10.1214/aos/1176325367.

  \bibitem{CaiLow2004} T. Tony Cai and Mark G. Low (2004), An Adaptation Theory for Nonparametric Confidence Intervals, Annals of Statistics 32(5), 1805--1840, DOI: 10.1214/009053604000000049.

  \bibitem{ArmstrongKolesar2018} Timothy B. Armstrong and Michal Kolesár (2018), Optimal Inference in a Class of Regression Models, Econometrica 86(2), 655--683, DOI: 10.3982/ECTA14434.

  \bibitem{GentzkowShapiroSinkinson2011} Matthew Gentzkow, Jesse M. Shapiro, and Michael Sinkinson (2011), The Effect of Newspaper Entry and Exit on Electoral Politics, American Economic Review 101(7), 2980--3018, DOI: 10.1257/aer.101.7.2980.

  \bibitem{HirshbergMalekiZubizarreta2021} David A. Hirshberg, Arian Maleki, and José R. Zubizarreta (2021), Minimax Linear Estimation of the Retargeted Mean, arXiv:1901.10296v2.

  \bibitem{AudibertTsybakov2011} Jean-Yves Audibert and Alexandre B. Tsybakov (2007), Fast Learning Rates for Plug-in Classifiers under the Margin Condition, arXiv:math/0507180v3.

  \bibitem{Gaiffas2005} Stéphane Gaïffas (2005), Convergence Rates for Pointwise Curve Estimation with a Degenerate Design, arXiv:math/0410354v3.

\end{thebibliography}

\end{document}
